%%%%%%%%%%%%%%%%%%%%%%%%%%%%%%%%%%%%%%%%%%%%%%%%%%%%%%%%%%%%%%%%%%%%%%
% Lớp 12 | Phương trình mặt phẳng, đường thẳng, mặt cầu | Phương trình đường thẳng trong không gian
% Dạng 2H5?2-3: Viết phương trình tổng quát, chính tắc, tham số đường thẳng
% Số câu: 81  (sinh tự động bởi tach_cau.py; từ nay thầy sửa trực tiếp file này)
%%%%%%%%%%%%%%%%%%%%%%%%%%%%%%%%%%%%%%%%%%%%%%%%%%%%%%%%%%%%%%%%%%%%%%

%%%%%%%%%%%%%%%% TRẮC NGHIỆM 4 PHƯƠNG ÁN %%%%%%%%%%%%%%%%

% nguồn: Phan1_Lop_12_Hinh_hoc_Chuong_5
\begin{ex}%[2H5N2-3]%[ma:5c122acb]
	Trong không gian $ Oxyz$, cho đường  thẳng $d$ có phương trình chính tắc $\dfrac{x-1}{3}=\dfrac{y+3}{2}=\dfrac{z-2}{-1}$. Phương trình tham số của đường thẳng $d$~là
	 \choice
		 {\True $\heva{x=&1+3t\\y=&-3+2t\\z=&2-t}$}
		 {$\heva{x=&-1+3t\\y=&3+2t\\z=&-2-t}$}
		{$\heva{x=&3-t\\y=&2+3t\\z=&-1-2t}$} 
		 {$\heva{x=&-3+t\\y=&-2-3t\\z=&1+2t}$}
	\loigiai{Phương trình tham số của đường thẳng $d$~là $\heva{x=&1+3t\\y=&-3+2t\\z=&2-t.}$ 
	 }
\end{ex}

% nguồn: Phan1_Lop_12_Hinh_hoc_Chuong_5
\begin{ex}%[2H5N2-3]%[ma:e836bffe]
	Trong  không gian $Oxyz$, phương trình đường thẳng đi qua điểm $P(1; 1;-1)$ và có vectơ chỉ phương $\overrightarrow{u}=(1; 2; 3)$ là
	 \choice
		{\True $\dfrac{x-1}{1}=\dfrac{y-1}{2}=\dfrac{z+1}{3}$} 
		 {$\dfrac{x+2}{1}=\dfrac{y+3}{2}=\dfrac{z+2}{3}$}
		 {$\dfrac{x-1}{1}=\dfrac{y-2}{1}=\dfrac{z-3}{-1}$}
		 {$\dfrac{x-1}{2}=\dfrac{y-1}{3}=\dfrac{z+1}{2}$}
	\loigiai{Phương  trình đường thẳng đi qua điểm $P(1; 1;-1)$ và có vectơ chỉ phương $\overrightarrow{u}=(1; 2; 3)$ là $\dfrac{x-1}{1}=\dfrac{y-1}{2}=\dfrac{z+1}{3}$.}
\end{ex}

% nguồn: Phan1_Lop_12_Hinh_hoc_Chuong_5
\begin{ex}%[2H5N2-3]%[ma:aee6bf85]
	Trong  không gian $Oxyz$, phương trình đường thẳng đi qua hai điểm $A(2;3;-1)$ và $B(2;3;-2)$~là
	\choice 
		{$\heva{& x=2+t\\ &y=3+t\\ &z=-1}$ $(t \in  \mathbb{R})$}
		{\True $\heva{&  x=2\\ &y=3\\ &z=-1-t}$ $(t \in \mathbb{R})$}
		{$\heva{& x=2 \\ &y=3+t \\ &z=-1+t}$ $(t \in \mathbb{R})$} 
		{$\heva{& x=2+t\\ &y=3\\  &z=-1+t}$ $(t \in \mathbb{R})$}
	\loigiai{ 
		Đường thẳng đi qua $A$  và nhận $\overrightarrow{AB}=(0;0;-1)$~làm véc-tơ chỉ phương nên có phương trình~là $\heva{& x=2\\ &y=3\\ &z=-1-t}\quad (t \in \mathbb{R})$.
	} 
\end{ex}

% nguồn: Phan1_Lop_12_Hinh_hoc_Chuong_5
\begin{ex}%[2H5N2-3]%[ma:bd355f5c]
	Trong không gian $Oxyz$, đường thẳng $\Delta$ đi qua $M(1;2;-3)$ và nhận vectơ $\vec{u}=(-1;2;1)$ làm vectơ chỉ phương có phương trình chính  tắc là
	\choice 
		{$\dfrac{x-1}{1}=\dfrac{y-2}{-2}=\dfrac{z+3}{1}$} 
		{$\dfrac{x+1}{-1}=\dfrac{y+2}{2}=\dfrac{z-3}{1}$} 
		{$\dfrac{x-1}{1}=\dfrac{y-2}{2}=\dfrac{z-3}{-1}$} 
		 {\True$\dfrac{x-1}{-1}=\dfrac{y-2}{2}=\dfrac{z+3}{1}$}
	 \loigiai{
		Phương trình đường thẳng  $\Delta \colon \dfrac{x-1}{-1}=\dfrac{y-2}{2}=\dfrac{z+3}{1}$.
	} 
\end{ex}

% nguồn: Phan1_Lop_12_Hinh_hoc_Chuong_5
\begin{ex}%[2H5N2-3]%[ma:9fdce70f]
	Trong không gian  $Oxyz$, phương trình đường thẳng đi qua điểm $M(1;-2;-5)$ và có vectơ chỉ phương $\overrightarrow{u}=(2;1;-1)$ là
	\choice 
		 {$\dfrac{x-1}{2}=\dfrac{y-2}{1}=\dfrac{z-5}{-1}$}
		{$\dfrac{x+1}{2}=\dfrac{y-2}{1}=\dfrac{z-5}{-1}$} 
		 {$\dfrac{x-2}{1}=\dfrac{y-1}{-2}=\dfrac{z+1}{5}$}
		 {\True $\dfrac{x-1}{2}=\dfrac{y+2}{1}=\dfrac{z+5}{-1}$}
	\loigiai{Phương trình đường thẳng đi qua điểm $M(1;-2;-5)$ và có vectơ chỉ  phương $\overrightarrow{u}=(2;1;-1)$ là $\dfrac{x-1}{2}=\dfrac{y+2}{1}=\dfrac{z+5}{-1}$.}
\end{ex}

% nguồn: Phan1_Lop_12_Hinh_hoc_Chuong_5
\begin{ex}%[2H5N2-3]%[ma:9802b06c]
	Trong không gian với hệ trục tọa độ $Oxyz$,  đường thẳng $d$ đi qua điểm $M(2;0;-1)$ và có vectơ chỉ phương $\overrightarrow{a}=(4;-6;2)$ là
	 \def\dotEX{}
	 \choice
		 {$\heva{&x=-2+2t \\&y=-3t \\&z=1+t.}$}
		 {$\heva{&x=-2+2t \\&y=-3t \\&z=-1+t.}$}
		{\True $\heva{&x=2+4t \\&y=-6t \\&z=-1+2t.}$} 
		 {$\heva{&x=-2+2t \\&y=6-3t \\&z=-3+t.}$}
	\loigiai{ 
		Phương trình đường thẳng $d$ là $\heva{&x=2+4t \\&y=-6t \\&z=-1+2t.}$ 
	 }
\end{ex}

% nguồn: Phan1_Lop_12_Hinh_hoc_Chuong_5
\begin{ex}%[2H5N2-3]%[ma:6524e9af]
	Trong không gian $Oxyz$, phương trình chính tắc  đường thẳng $d$ đi qua $M(3;2;1)$ và nhận vectơ $\overrightarrow{u} = (1;2;-3)$ làm một vectơ chỉ phương là
	 \choice
		 {$\dfrac{x-1}{3}=\dfrac{y-2}{-2}=\dfrac{z+3}{1}$}
		{\True  $\dfrac{x-3}{1}=\dfrac{y-2}{2}=\dfrac{z-1}{-3}$}
		{$\dfrac{x-3}{1}=\dfrac{y+2}{2}=\dfrac{z-1}{-3}$} 
		 {$\dfrac{x-1}{3}=\dfrac{y-2}{2}=\dfrac{z+3}{1}$}
	\loigiai{ 
		phương trình chính tắc đường thẳng $d$ đi qua $M(3;2;1)$ và nhận  vectơ $\overrightarrow{u} = (1;2;-3)$ làm một vectơ chỉ phương là
		 $$\dfrac{x-3}{1}=\dfrac{y-2}{2}=\dfrac{z-1}{-3}.$$
	} 
\end{ex}

% nguồn: Phan1_Lop_12_Hinh_hoc_Chuong_5
\begin{ex}%[2H5N2-3]%[ma:54133705]
	Trong không gian với hệ tọa~độ $  Oxyz$, đường thẳng nhận ${\vec u}\left(a;b;c\right)\left({abc}\ne 0\right)$~là véc-tơ chỉ phương và đi qua điểm ${I}\left({x}_0;{{y}_0};{{z}_0}\right)$ có phương trình chính tắc~là
	 \choice
		{$\dfrac{x-x_0}{a}=\dfrac{y-y_0}{c}=\dfrac{z-z_0}{b}$} 
		 {\True $\dfrac{x-x_0}{a}=\dfrac{y-y_0}{b}=\dfrac{z-z_0}{c}$}
		{$\dfrac{x-x_0}{a}=\dfrac{y-y_0}{-b}=\dfrac{z-z_0}{c}$} 
		{$\dfrac{x-x_0}{a}=\dfrac{y-y_0}{b}=\dfrac{z-z_0}{-c}$} 
	\loigiai{Đường thẳng nhận ${\vec u}\left(a;b;c\right)\left({abc}\ne 0\right)$~là véc-tơ chỉ phương và đi qua điểm ${I}\left({x}_0;{{y}_0};{{z}_0}\right)$ có phương trình chính tắc~là  $\dfrac{x-x_0}{a}=\dfrac{y-y_0}{b}=\dfrac{z-z_0}{c}$.
	 }
\end{ex}

% nguồn: Phan1_Lop_12_Hinh_hoc_Chuong_5
\begin{ex}%[2H5N2-3]%[ma:baf4414d]
	Trong không gian $Oxyz$, phương trình tham số của đường thẳng đi qua điểm $M(2;0;-1)$ và có vectơ chỉ phương  $\vec{a}=(2;-3;1)$ là
	 \choice
		 {$\heva{&x=2+2t\\&y=-3\\&z=1-t}$}
		 {$\heva{&x=-2+4t\\&y=-6t\\&z=1+2t}$}
		 {\True$\heva{&x=2+2t\\&y=-3t\\&z=-1+1t}$}
		{$\heva{&x=-2+2t\\&y=-3t\\&z=1+t}$} 
	\loigiai{ 
		Đường thẳng đi qua $M(2;0;-1)$ có vectơ chỉ  phương $\vec{a}=(2;-3;1)$ có dạng $\heva{&x=2+2t\\&y=-3t\\&z=-1+t}$.
	 }
\end{ex}

% nguồn: Phan1_Lop_12_Hinh_hoc_Chuong_5
\begin{ex}%[2H5N2-3]%[ma:3ef3a981]
	Trong không gian $ Oxyz$, cho hai điểm $ A\left(1;-2;-3\right),B\left(2;4;5\right)$. Phương trình chính tắc của đường thẳng $ AB$~là 
	 \choice
		{$\dfrac{x+1}{1}=\dfrac{y-2}{6}=\dfrac{z-3}{8}$} 
		{$\dfrac{x+1}{-1}=\dfrac{y+2}{6}=\dfrac{z+3}{8}$} 
		{$\dfrac{x-1}{1}=\dfrac{y+2}{-6}=\dfrac{z+3}{8}$} 
		 {\True $\dfrac{x-1}{1}=\dfrac{y+2}{6}=\dfrac{z+3}{8}$}
	\loigiai{ 
		Phương trình đường  thẳng $ AB$ có một véc-tơ chỉ phương $\vec{u}=\overrightarrow{AB}=(1;6;8)$. Phương trình chính tắc của đường thẳng $AB$~là $\dfrac{x-1}{1}=\dfrac{y+2}{6}=\dfrac{z+3}{8}$.
	 }
\end{ex}

% nguồn: Phan1_Lop_12_Hinh_hoc_Chuong_5
\begin{ex}%[2H5N2-3]%[ma:41c7cf72]
	Trong không gian với hệ tọa  độ $Oxyz$, cho hai điểm $M(1; 0; -1)$ và $N(3; 2; -9)$. Đường thẳng $MN$ có phương trình là
	 \choice
		{$\dfrac{x+1}{1} = \dfrac{y}{1} = \dfrac{z-1}{-4}$} 
		{$\dfrac{x-1}{1} = \dfrac{y}{1} =  \dfrac{z+1}{4}$}
		{$\dfrac{x+1}{1} = \dfrac{y}{1}  = \dfrac{z-1}{4}$}
		 {\True $\dfrac{x-1}{1} = \dfrac{y}{1} = \dfrac{z+1}{-4}$}
	\loigiai{ 
		Ta có $\overrightarrow{MN}=  (2; 2; -8)=2 (1; 1; -4)$.\\
		Đường thẳng $MN$ đi qua điểm $M(1; 0; -1)$ và có một véctơ chỉ phương  là $\overrightarrow{u}=(1;1;-4)$.\\
		Vậy phương trình đường thẳng $MN$  là
		$$ \dfrac{x-1}{1} = \dfrac{y}{1} = \dfrac{z+1}{-4}. $$ 
	 }
\end{ex}

% nguồn: Phan1_Lop_12_Hinh_hoc_Chuong_5
\begin{ex}%[2H5N2-3]%[ma:a36e77f2]
	Trong không gian $ Oxyz$, đường thẳng đi qua điểm $ A\left(-12;13;-14\right)$ và có véc-tơ chỉ phương $\overrightarrow u\left(17;-18;-19\right)$ có phương trình  chính tắc~là
	\choice 
		 {$\dfrac{x+12}{17}=\dfrac{y-13}{18}=\dfrac{z+14}{-19}$}
		{$\dfrac{x+12}{17}=\dfrac{y-13}{-18}=\dfrac{z+14}{19}$} 
		 {$\dfrac{x+12}{17}=\dfrac{y-13}{18}=\dfrac{x+14}{19}$}
		{\True  $\dfrac{x+12}{17}=\dfrac{y-13}{-18}=\dfrac{z+14}{-19}$}
	\loigiai{ 
		Đường thẳng  đi qua điểm $ A\left(-12;13;-14\right)$ và có véc-tơ chỉ phương $\overrightarrow u\left(17;-18;-19\right)$ có phương trình chính tắc~là $\dfrac{x+12}{17}=\dfrac{y-13}{-18}=\dfrac{z+14}{-19}$.
	 }
\end{ex}

% nguồn: Phan1_Lop_12_Hinh_hoc_Chuong_5
\begin{ex}%[2H5N2-3]%[ma:8d2b071b]
	Trong không  gian với hệ tọa độ $Oxyz$, phương trình nào sau đây là phương trình chính tắc của một đường thẳng?
	 \choice
		{\True $\dfrac{x-9}{2}=\dfrac{y-10}{3}=\dfrac{z-11}{4}$} 
		{$\dfrac{x^2-9}{2}=\dfrac{y-10}{3}=\dfrac{z-11}{4}$} 
		 {$\dfrac{x-9}{2}=\dfrac{y^2-10}{3}=\dfrac{z-11}{4}$}
		 {$\dfrac{x-9}{2}=\dfrac{y-10}{3}=\dfrac{z^2-11}{4}$}
	\loigiai{Phương trình chính tắc của  một đường thẳng là phương trình $\dfrac{x-9}{2}=\dfrac{y-10}{3}=\dfrac{z-11}{4}$.
	 }
\end{ex}

% nguồn: Phan1_Lop_12_Hinh_hoc_Chuong_5
\begin{ex}%[2H5N2-3]%[ma:e629a173]
	Trong không gian với hệ tọa độ $Oxyz$, cho đường thẳng $d$ đi qua điểm $M(0;-1; 4)$ và nhận vectơ $\overrightarrow{u}=(3;-1; 5)$ làm vectơ chỉ phương. Phương trình tham số  của $d$ là
	 \choice
		 {$\heva{&x=3t \\
		 &y=1-t \\
		&z=4+5t 
		 &}$}
		 {$\heva{&x=3\\
		&y=-1-t  \\
		 &z=5+4t
		&}$} 
		 {\True $\heva{&x=3t \\
		 &y=-1-t \\
		 &z=4+5t
		 &}$}
		 {$\heva{&x=3t \\
		&y=1-t \\ 
		 &z=-4+5t
		&}$} 
	\loigiai{ 
		Đường thẳng $d$ đi qua điểm  $M(0;-1; 4)$ và nhận vectơ $\overrightarrow{u}=(3;-1; 5)$ làm vectơ chỉ phương. Phương trình tham số của $d$ là $\heva{&x=3t \\&y=-1-t \\&z=4+5t.}$
	} 
\end{ex}

% nguồn: Phan1_Lop_12_Hinh_hoc_Chuong_5
\begin{ex}%[2H5N2-3]%[ma:af729e39]
	Trong không gian $Oxyz$, phương trình  đường thẳng $\Delta$ đi qua điểm $M(1;-2;4)$ và có véc-tơ chỉ phương $\vec{u}=(2;3;-1)$~là
	 \choice
		{\True $\heva{& x=1+2t\\ &y=-2+3t\\ &z=4-t}$ $(t  \in \mathbb{R})$}
		{$\heva{&  x=1+2t\\ &y=-2+3t\\ &z=4+t}$ $(t\in\mathbb{R})$}
		{$\heva{& x=2+t\\  &y=3-2t\\ &z=-1+4t}$ $(t\in\mathbb{R})$}
		{$\heva{&  x=2+t\\ &y=3+2t\\ &z=-1+4t}$ $(t\in\mathbb{R})$}
	 \loigiai{
		Đường  thẳng $\Delta$ đi qua điểm $M(1;-2;4)$ và có véc-tơ chỉ phương $\vec{u}=(2;3;-1)$ có phương trình
		 là $\heva{& x=1+2t\\ &y=-2+3t\\ &z=4-t}$ $(t \in \mathbb{R})$.
	 }
\end{ex}

% nguồn: Phan1_Lop_12_Hinh_hoc_Chuong_5
\begin{ex}%[2H5N2-3]%[ma:e144b9ca]
	Trong không gian $Oxyz$, phương trình của mặt phẳng đi qua điểm $K(1;0;-2)$ và có một véc-tơ pháp tuyến $\overrightarrow{n}=(1;3;-5)$  là
	\choice 
		{$x-2z+9=0$} 
		{$x+3y-5z+9=0$} 
		 {$x-2z-11=0$}
		{\True $x+3y-5z-11=0$} 
	\loigiai{ 
		Phương trình của mặt  phẳng đi qua điểm $K(1;0;-2)$ và có một véc-tơ pháp tuyến $\overrightarrow{n}=(1;3;-5)$ là
		 \[1(x-1)+3(y-0)+(-5)(z+2)=0\Leftrightarrow x+3y-5z-11=0.\]
		Phương trình mặt phẳng cần tìm  là $x+3y-5z-11=0$.
	 }
\end{ex}

% nguồn: Phan1_Lop_12_Hinh_hoc_Chuong_5
\begin{ex}%[2H5N2-3]%[ma:17f6664b]
	Trong không gian $Oxyz$, cho đường thẳng $d$ đi qua điểm $M(2;2;1)$ và có một vectơ chỉ phương $\overrightarrow{n}=(2;3;-1)$.  Phương trình của $d$ là
	\choice 
		{\True $\heva{& x=2+5t \\ & y=2+2t \\ & -1-3t}$} 
		 {$\heva{& x=2+5t \\ & y=2+2t \\ & 1+3t}$}
		 {$\heva{& x=2+5t \\ & y=2+2t \\ & 1-3t}$}
		{$\heva{& x=5+2t \\ & y=2+2t \\ & -3-1t}$} 
	\loigiai{ 
	Phương trình đường thẳng đi  qua điểm $M(2;2;1)$ và có một vectơ chỉ phương $\overrightarrow{n}=(2;3;-1)$ là $d\colon\heva{& x=2+5t \\ & y=2+2t \\ & -1-3t.}$
	 }
\end{ex}

% nguồn: Phan1_Lop_12_Hinh_hoc_Chuong_5
\begin{ex}%[2H5N2-3]%[ma:8f07717b]
	Trong không gian với hệ tọa~độ  $Oxyz$, đường thẳng nhận ${\vec u}\left(a;b;c\right)$~là véc-tơ chỉ phương và đi qua điểm ${I}\left({x}_0;{{y}_0};{{z}_0}\right)$ có phương trình tham số~là
	 \choice
		{$\heva{x=&a+x_0t\\y=&b+y_0t\\z=&c+z_0t}$} 
		{$\heva{x=&x_0+at\\y=&y_0+ct\\z=&z_0+bt}$} 
		{$\heva{x=&x_0+at\\y=&y_0+bt\\z=&z_0-ct}$} 
		{\True  $\heva{x=&x_0+at\\y=&y_0+bt\\z=&z_0+ct}$}
	\loigiai{Đường thẳng nhận ${\vec u}\left(a;b;c\right)$~là véc-tơ chỉ phương và đi qua điểm ${I}\left({x}_0;{{y}_0};{{z}_0}\right)$ có  phương trình tham số~là$\heva{x=&x_0+at\\y=&y_0+bt\\z=&z_0+ct}$.
	 }
\end{ex}

% nguồn: Phan1_Lop_12_Hinh_hoc_Chuong_5
\begin{ex}%[2H5N2-3]%[ma:086634de]
	Trong không gian $Oxyz$, đường thẳng đi qua điểm $B(-1; 3; 6)$ và nhận $\vec{u}=(2;-3; 8)$ làm vectơ chỉ phương có  phương trình chính tắc là
	 \choice
		{$\dfrac{x+1}{2}=\dfrac{y-3}{3}=\dfrac{z-6}{8}$} 
		{\True $\dfrac{x+1}{2}=\dfrac{y-3}{-3}=\dfrac{z-6}{8}$} 
		 {$\dfrac{x-1}{2}=\dfrac{y+3}{-3}=\dfrac{z+6}{8}$}
		{$\dfrac{x+1}{-2}=\dfrac{y-3}{3}=\dfrac{z-6}{8}$} 
	 \loigiai{
		Đường thẳng đi qua $B$ và nhận  $\vec{u}=(2;-3; 8)$ làm vectơ chỉ phương có phương trình chính tắc là
		 $$\dfrac{x+1}{2}=\dfrac{y-3}{-3}=\dfrac{z-6}{8}.$$
	} 
\end{ex}

% nguồn: Phan1_Lop_12_Hinh_hoc_Chuong_5
\begin{ex}%[2H5N2-3]%[ma:bc3e733f]
	Trong không gian $Oxyz$, phương  trình tham số của đường thẳng đi qua điểm $A(0;8;3)$ và có một véc-tơ chỉ phương $\overrightarrow{u}=(2;1;-1)$ là
	\choice 
		{$\heva{&x=2\\&y=8+t\\&z=-1+3t},t\in\mathbb{R}$} 
		{$\heva{&x=2t\\&y=8+t\\&z=3+t},t\in\mathbb{R}$} 
		{$\heva{&x=2t\\&y=8\\&z=3-t},t\in\mathbb{R}$} 
		{\True  $\heva{&x=2t\\&y=8+t\\&z=3-t},t\in\mathbb{R}$}
	 \loigiai{
		Phương trình đường thẳng $d$ đi qua điểm $A(0;8;3)$ và có  véc-tơ chỉ phương $\overrightarrow{u}=(2;1;-1)$ là
		\[\heva{&x=2t\\&y=8+t\\&z=3-t},t\in\mathbb{R}.\] 
	} 
\end{ex}

% nguồn: Phan1_Lop_12_Hinh_hoc_Chuong_5
\begin{ex}%[2H5N2-3]%[ma:69ae0ece]
	Trong không gian $Oxyz$, cho đường thẳng $\Delta\colon  \heva{& x=1+t\\ &y=-2+2t\\ &z=3-t}$ $(t \in \mathbb{R})$. Điểm nào sau đây thuộc đường thẳng $\Delta$?
	 \choice
		 {$M(1;2;-1)$}
		 {\True $N(2;0;2)$}
		 {$P(-1;2;5)$}
		{$Q(1;-2;2)$} 
	\loigiai{ 
		Thay lần lượt tọa~độ của các điểm $M$, $N$, $P$, $Q$ vào phương trình của đường thẳng $\Delta$, ta thấy tọa~độ của điểm $N$ thoả mãn nên điểm $N$ thuộc  đường thẳng $\Delta$.
	 }
\end{ex}

% nguồn: Phan1_Lop_12_Hinh_hoc_Chuong_5
\begin{ex}%[2H5N2-3]%[ma:a12c33db]
	Trong không gian $Oxyz$, đường thẳng $d$ đi qua điểm $M(3;-1;4)$ và có vectơ chỉ phương $\overrightarrow{u}  = (-2;4;5)$ có phương trình tham số là
	\choice 
		 {\True $\heva{&x=3-2t\\&y=-1+4t\\&z=4+5t}$}
		{$\heva{&x=3-2t\\&y=1+4t\\&z=4+5t}$} 
		{$\heva{&x=3+2t\\&y=-1+4t\\&z=4+5t}$} 
		{$\heva{&x=-2+3t\\&y=4-t\\&z=5+4t}$} 
	\loigiai{Đường thẳng đi qua điểm $M(3;-1;4)$ và có vectơ chỉ phương $\overrightarrow{u}=(-2;4;5)$  có phương trình tham số là $\heva{&x=3-2t\\&y=-1+4t\\&z=4+5t.}$
	} 
\end{ex}

% nguồn: Phan1_Lop_12_Hinh_hoc_Chuong_5
\begin{ex}%[2H5N2-3]%[ma:a95fb29e]
	[GV:Xuan Vy  Pham]
	Phương trình đường thẳng đi qua $A(1;-2;0)$ và vuông góc với  mặt phẳng \\$(P) \colon x-2y+2z+1=0$ là
	\choice 
		{\True $\dfrac{x-1}{1}=\dfrac{y+2}{-2}=\dfrac{z}{2}$} 
		 {$\dfrac{x-1}{1}=\dfrac{y+2}{2}=\dfrac{z}{2}$}
		 {$\dfrac{x-1}{-1}=\dfrac{y-2}{-2}=\dfrac{z}{-2}$}
		{$\dfrac{x-1}{1}=\dfrac{y-2}{-2}=\dfrac{z}{2}$} 
	\loigiai{Đường thẳng $d$ vuông  góc với mặt phẳng $(P) \colon x-2y+2z+1=0$ nên một vectơ chỉ phương của đường thẳng $d$ là vectơ pháp tuyến $\overrightarrow{n_P}=(1;-2;2)$ của mặt phẳng $(P)$.\\
	Phương trình  đường thẳng $d$ đi qua $A(1;-2;0)$ có vectơ chỉ phương $\overrightarrow{n_P}=(1;-2;2)$ là $$\dfrac{x-1}{1}=\dfrac{y+2}{-2}=\dfrac{z}{2}.$$}
\end{ex}

% nguồn: Phan1_Lop_12_Hinh_hoc_Chuong_5
\begin{ex}%[2H5N2-3]%[ma:ac32773e]
	Trong không gian  $Oxyz$, cho đường thẳng $d$ có phương trình $\heva{&x=1+3t\\&y=8-t\\&z=6t}$ trong đó $t\in\mathbb{R}$. Một véc-tơ chỉ phương của $d$ là
	 \choice
		 {\True $\overrightarrow{u}_{2}=(3;-1;6)$}
		{$\overrightarrow{u}_{1}=(1;8;6)$} 
		 {$\overrightarrow{u}_{3}=(1;8;0)$}
		{$\overrightarrow{u}_{4}=(3;-1;0)$} 
	\loigiai{ 
		Một véc-tơ chỉ phương của đường thẳng  $d$ là $\overrightarrow{u}_{2}=(3;-1;6)$.
	} 
\end{ex}

% nguồn: Phan1_Lop_12_Hinh_hoc_Chuong_5
\begin{ex}%[2H5N2-3]%[ma:d471bb53]
	Trong không gian $ Oxyz$, đường thẳng đi qua điểm ${A}\left(3;2;3\right)$  và có véc-tơ chỉ phương ${\vec u}\left(1;-2;1\right)$ có phương trình tham số~là
	\renewcommand{\dotEX}{} 
	 \choice
		{$\heva{x=&3-t\\y=&2-2t\\z=&3+t.}$} 
		 {\True $\heva{x=&3+t\\y=&2-2t\\z=&3+t.}$}
		{$\heva{x=&3+t\\y=&2-2t\\z=&3-t.}$} 
		{$\heva{x=&-3+t\\y=&2-2t\\z=&3+t.}$} 
	 \loigiai{
		Đường thẳng đi qua điểm ${A}\left(3;2;3\right)$ và có  véc-tơ chỉ phương ${\vec u}\left(1;-2;1\right)$ có phương trình tham số~là $\heva{x=&3+t\\y=&2-2t\\z=&3+t.}$
	} 
\end{ex}

% nguồn: Phan1_Lop_12_Hinh_hoc_Chuong_5
\begin{ex}%[2H5N2-3]%[ma:8c6a86ac]
	Trong không  gian $Oxyz$, đường thẳng $d$ đi qua điểm $M(1;1;1)$ có vec-tơ chỉ phương $\vec{u} = (1;2;3)$ có phương trình là
	\choice 
		{$\heva{&x=1+t \\&y = 1+2t  \\&z=1-3t}(t\in\mathbb{R})$}
		{$\heva{&x=1+3t \\&y = 1+ 2t \\&z=1+t}(t\in  \mathbb{R})$}
		{$\heva{&x=1+t \\&y=1+3t \\&z=1+2t}(t\in\mathbb{R})$} 
		{\True $\heva{&x=1+t \\&y = 1+2t \\&z=1+3t}(t\in \mathbb{R})$} 
	 \loigiai{
		Đường thẳng $d$ đi qua điểm $M(1;1;1)$  có vectơ chỉ phương $\vec{u} = (1;2;3)$ có phương trình là\\
		$\heva{&x=1+t \\&y  = 1+2t \\&z=1+3t}(t\in \mathbb{R})$
	} 
\end{ex}

% nguồn: Phan1_Lop_12_Hinh_hoc_Chuong_5
\begin{ex}%[2H5N2-3]%[ma:0e1eb6ef]
	Trong không gian $Oxyz$, phương trình đường thẳng $d$ đi qua điểm $M(2;1;-1)$ và có  một vectơ chỉ phương $\overrightarrow{u} = (1;-2;3)$ là
	 \choice
		{\True  $\dfrac{x - 2}{1} = \dfrac{y - 1}{-2} = \dfrac{z + 1}{3}$}
		{$\dfrac{x + 1}{2} = \dfrac{y - 2}{1} = \dfrac{z + 3}{-1}$} 
		{$\dfrac{x - 1}{2} = \dfrac{y  + 1}{-2} = \dfrac{z - 1}{3}$}
		{$\dfrac{x - 1}{2} = \dfrac{y + 2}{1}  = \dfrac{z - 3}{-1}$}
	 \loigiai{
		Đường thẳng đi  qua $M(2;1;-1)$, nhận $\overrightarrow{u} = (1;-2;3)$ làm vectơ chỉ phương có phương trình là
		$$ 
		\dfrac{x - 2}{1} =  \dfrac{y - 1}{-2} = \dfrac{z + 1}{3}.
		 $$
	} 
\end{ex}

% nguồn: Phan1_Lop_12_Hinh_hoc_Chuong_5
\begin{ex}%[2H5N2-3]%[ma:f72d9888]
	Trong không gian $Oxyz$,  phương trình chính tắc của đường thẳng $d$ đi qua hai điểm $A(1;3;2)$, $B(2;1;1)$ là
	\choice 
		{\True $\dfrac{x-1}{1}=\dfrac{y-3}{-2}=\dfrac{z-2}{-1}$} 
		{$\dfrac{x-1}{1}=\dfrac{y+3}{-2}=\dfrac{z-2}{-1}$} 
		{$\dfrac{x-1}{1}=\dfrac{y-3}{2}=\dfrac{z-2}{-1}$} 
		 {$\dfrac{x-1}{1}=\dfrac{y+3}{-1}=\dfrac{z-2}{-1}$}
	 \loigiai{
		Gọi $\overrightarrow{AB}=(1;-2;-1)$ là  vectơ chỉ phương của đường thẳng $d$.\\
		Phương trình chính tắc của đường  thẳng đi qua $A(1;3;2)$ và có vectơ chỉ phương $(1;-2;-1)$ là
		\[ 
		d\colon \dfrac{x-1}{1}=\dfrac{y-3}{-2}=\dfrac{z-2}{-1}. 
		\] 
	 }
\end{ex}

% nguồn: Phan1_Lop_12_Hinh_hoc_Chuong_5
\begin{ex}%[2H5N2-3]%[ma:a55ceba7]
	Trong không  gian $ Oxyz$, cho ba điểm $ A\left(1;-1;1\right)$, $B\left(-1;0;2\right)$, $C\left(2;1;3\right)$, đường thẳng qua $ A$ và song song với đường thẳng $BC$ có phương trình~là
	\renewcommand{\dotEX}{} 
	 \choice
		 {$\heva{x=&1-3t\\y=&-1+t\\z=&1+t.}$}
		{$\heva{x=&1+3t\\y=&-1-t\\z=&1+t.}$} 
		{$\heva{x=&1+3t\\y=&-1-t\\z=&1-t.}$} 
		{\True $\heva{x=&1+3t\\y=&-1+t\\z=&1+t.}$} 
	\loigiai{ 
		Đường thẳng $d$ qua $ A$ và  song song với đường thẳng $BC$ có véc-tơ chỉ phương~là $\vec{u}=\overrightarrow{BC}=(3;1;1)$. Vậy phương trình đường thẳng $d$~là $\heva{x=&1+3t\\y=&-1+t\\z=&1+t.}$
	} 
\end{ex}

% nguồn: Phan1_Lop_12_Hinh_hoc_Chuong_5
\begin{ex}%[2H5N2-3]%[ma:536d55bb]
	Trong không gian $Oxyz$, phương trình tham số của đường thẳng $(d)$ đi qua điểm $M(1;-1;2)$  và có vectơ chỉ phương $\overrightarrow u=(2;4;-5)$ là
	 \choice
		{$\heva{&x=-1+2t\\& y=1+4t\\&z=-2-5t}$} 
		{\True  $\heva{&x=1+2t\\&y=-1+4t\\&z=2-5t}$}
		 {$\heva{&x=2+t\\&y=4-t\\&z=-5+2t}$}
		 {$\heva{&x=-2+t\\&y=-4-t\\&z=5+2t}$}
	\loigiai{ 
		Phương  trình tham số của đường thẳng $(d)$ đi qua điểm $M(1;-1;2)$ và có vectơ chỉ phương $\overrightarrow u=(2;4;-5)$ là
		$$\heva{&x=1+2t\\& y=-1+4t\\&  z=2-5t}\ (t\in\mathbb{R}).$$
	 }
\end{ex}

% nguồn: Phan1_Lop_12_Hinh_hoc_Chuong_5
\begin{ex}%[2H5H2-3]%[ma:48c081d3]
	Trong không gian với hệ tọa độ $Oxyz$, cho điểm $P(2;1;-4)$ và đường thẳng $d_1 \colon \heva{&x=t\\&y=1-t\\ &z=-12+t}$. Đường thẳng $\Delta$ cắt $d_1$ và trục $Oy$ lần lượt tại $A$ và $B$ sao cho $P$ là trung điểm của $AB$. Đường thẳng $\Delta$ có phương  trình là
	\choice 
		{$\heva{&x=2+3t\\ &y=1+2t\\  &z=-4+t}$}
		{\True $\heva{&x=2t\\ &y=-2+3t\\ &z=-4t}$} 
		{$\heva{&x=2+2t\\  &y=1-5t\\&z=-4+t}$}
		{$\heva{&x=3-t\\ &y=3-2t\\  &z=1-5t}$}
	 \loigiai{
		Do $B=\Delta\cap Oy$ suy ra $B\in Oy$  nên $B(0;a;0)$.\\
		Vì $P$ là trung điểm $AB$, suy ra  $A(4;2-a;-8)$.\\
		Theo giả thiết, $A\in d_{1}$ nên  $\heva{&4=t\\ &2-a=t\\ &-8=-12+t} \Leftrightarrow \heva{&t=4\\&a=-2}$.\\
		 Suy ra $A(4;4;-8)$ và $B(0;-2;0)$.\\
		Đường thẳng $\Delta$ đi qua điểm  $B(0;-2;0)$ và nhận $\vec{PA}=(2;3;-4)$ làm vectơ chỉ phương có phương trình là $\Delta \colon \heva{&x=2t\\ &y=-2+3t\\ &z=-4t}$.
	 }
\end{ex}

% nguồn: Phan1_Lop_12_Hinh_hoc_Chuong_5
\begin{ex}%[2H5H2-3]%[ma:bd9a5560]
	Trong không gian $Oxyz$,  phương trình của đường thẳng $\Delta$ đi qua điểm $M(2;0;-1)$ và có vectơ chỉ phương $\overrightarrow{a}=(2;-3;1)$ là
	\choice 
		{\True  $\dfrac{x-2}{2}=\dfrac{y}{-3}=\dfrac{z+1}{1}$}
		 {$\dfrac{x+2}{2}=\dfrac{y}{-3}=\dfrac{z-1}{1}$}
		{$\dfrac{x-2}{2}=\dfrac{y}{3}=\dfrac{z+1}{1}$} 
		{$\dfrac{x+2}{2}=\dfrac{y}{3}=\dfrac{z-1}{1}$} 
	 \loigiai{
		Phương trình của đường thẳng $\Delta$ đi qua điểm $M(2;0;-1)$ và có vectơ chỉ phương $\overrightarrow{a}=(2;-3;1)$ là\\ 
		 $$\dfrac{x-2}{2}=\dfrac{y}{-3}=\dfrac{z+1}{1}.$$
	} 
\end{ex}

% nguồn: Phan1_Lop_12_Hinh_hoc_Chuong_5
\begin{ex}%[2H5H2-3]%[ma:5b383c88]
	Trong không gian $Oxyz$, cho điểm $M(3;-1;2)$  và đường thẳng $\Delta:\dfrac{x-1}{2}=\dfrac{y-5}{-5}=\dfrac{z+2}{-4}$. Gọi $d$ là đường thẳng đi qua $M$ và song song với đường thẳng $\Delta$. Phương trình nào sau đây là phương trình tham số của đường thẳng $d$?
	 \choice
		{\True $\left\{\begin{aligned}x &= 3+2t \\ y &= -1-5t \\ z &= 2-4t
\end{aligned}\right.$}
{$\left\{\begin{aligned}x &= 2+3t \\ y &= -5-t \\ z &= -4+2t
\end{aligned}\right.$}
{$\left\{\begin{aligned}x &= 3+t \\ y &= -1+5t \\ z &= 2-2t
\end{aligned}\right.$}
{$\left\{\begin{aligned}x &= 3+2t \\ y &= -1+5t \\ z &= 2-4t
\end{aligned}\right.$}
\loigiai{ 
	$\Delta$ có một vectơ chỉ phương $\vec{u}=(2;-5;-4)$. \\ 
	Vì $d$ song song với $\Delta$  nên $\vec{u}$ cũng là một vectơ chỉ phương của $d$. \\
	Phương trình tham số của $d$ là $\left\{\begin{aligned}x &= 3+2t \\ y &= -1-5t \\ z &= 2-4t
	\end{aligned}\right.$.
	 }
	\end{ex}

% nguồn: Phan1_Lop_12_Hinh_hoc_Chuong_5
\begin{ex}%[2H5H2-3]%[ma:f4755396]
	Trong không gian với hệ tọa độ $Oxyz$, cho điểm $N(0;-1;2)$ và hai đường thẳng $d_1 \colon \heva{&x=1-t\\&y=2+t\\& z=3-2t}$ và  $d_2 \colon \dfrac{x-2}{1} = \dfrac{y+1}{-2} = \dfrac{z-11}{4}.$ Đường thẳng $\Delta$ cắt $d_1$, $d_2$ lần lượt tại $A$ và $B$ sao cho $N$ là trung điểm của $AB$. Đường thẳng $\Delta$ có phương trình là
	\choice 
		{$\heva{&x=0\\  &y=-1+t\\ &z=2+3t}$}
		{$\heva{&x=0\\ &y=1+4t\\ &z=-1-6t}$} 
		 {$\heva{&x=0\\ &y=-1+5t\\ &z=2+3t}$}
		{\True $\heva{&x=0\\ &y=3+4t\\ &z=1-t}$} 
	 \loigiai{
		Do $A=\Delta\cap d_{1}$ suy  ra $A(1-t;2+t;3-2t)$.\\
		Vì $N$ là trung điểm  $AB$, suy ra $B(t-1;-4-t;2t+1)$.\\
		Theo giả  thiết, $B\in d_{2}$ nên $\dfrac{t-1-2}{1}=\dfrac{-4-t+1}{-2}=\dfrac{2t+1-11}{4}\Leftrightarrow t=1$.\\
		Suy ra $A(0;3;1)$,  $B(0;-5;3)$ và $\vec{AB}=(0;-8;2)$.\\
		Đường thẳng $\Delta$ đi qua hai điểm $A(0;3;1)$ và nhận $\vec{u}=-\dfrac{1}{2}\vec{AB}=(0;4;-1)$ làm vectơ  chỉ phương có phương trình là $\Delta\colon\heva{&x=0\\ &y=3+4t\\ &z=1-t}$.
	 }
\end{ex}

% nguồn: Phan1_Lop_12_Hinh_hoc_Chuong_5
\begin{ex}%[2H5H2-3]%[ma:7c7904c3]
	Trong không gian $Oxyz$, đường thẳng đi qua điểm $A(1;2;3)$ và  song song với đường thẳng $d\colon \dfrac{x}{-1}=\dfrac{y}{3}=\dfrac{z}{-2}$ có phương trình chính tắc~là
	\choice 
		 {$\dfrac{x-1}{1}=\dfrac{y+3}{2}=\dfrac{z-2}{3}$}
		{\True  $\dfrac{x-1}{1}=\dfrac{y-2}{-3}=\dfrac{z-3}{2}$}
		 {$\dfrac{x+1}{1}=\dfrac{y-3}{2}=\dfrac{z+2}{3}$}
		 {$\dfrac{x+1}{-1}=\dfrac{y+2}{3}=\dfrac{z+3}{-2}$}
	 \loigiai{
		Đường thẳng $d$ có  một véc-tơ chỉ phương~là $\vec{u}=(-1;3;-2)$.\\
		Đường thẳng $\Delta$ song song với $d$ nên có véc-tơ chỉ phương~là $\vec{a}=(1;-3;2)$.\\ 
		Vậy đường  thẳng $\Delta$ có phương trình chính tắc~là $\dfrac{x-1}{1}=\dfrac{y-2}{-3}=\dfrac{z-3}{2}$.
	} 
\end{ex}

% nguồn: Phan1_Lop_12_Hinh_hoc_Chuong_5
\begin{ex}%[2H5H2-3]%[ma:910e88d1]
	Trong không  gian $Oxyz$, đường thẳng $d$ đi qua hai điểm $M(1;-1;3)$, $N(3;0;2)$ có phương trình tham số là
	\choice 
		 {\True $\heva{&
		x=3+2t \\& 
		 y=t \\&
		 z=2-t}$}
		{$\heva{& 
		x=1+2t  \\&
		y=-1+t \\& 
		z=3+t}$} 
		 {$\heva{&
		x=2+t \\& 
		y=1-t \\& 
		z=-1+3t}$} 
		{$\heva{& 
		x=1+2t \\& 
		y=1+t \\& 
		 z=3-t}$}
	\loigiai{Đường thẳng  $d$ đi qua hai điểm $M(1;-1;3)$, $N(3;0;2)$ nên có vectơ chỉ phương $\overrightarrow{MN}=(2;1;-1)$.\\
		Vậy $d$ có  phương trình tham số là $\heva{&
		 x=3+2t \\&
		 y=t \\&
	z=2-t.}$} 
\end{ex}

% nguồn: Phan1_Lop_12_Hinh_hoc_Chuong_5
\begin{ex}%[2H5H2-3]%[ma:e0bd83f4]
	Trong không gian $Oxyz$, cho đường thẳng $\Delta\colon \dfrac{x-8}{2}=\dfrac{y+2}{4}=\dfrac{z-3}{m-1}$ $(m \neq 1)$ và điểm $M(6;-6;6)$ sao  cho $M$ thuộc $\Delta$. Giá trị của $m$ bằng
	 \choice
		 {\True $-2$}
		 {$2$}
		{$4$} 
		{$-4$} 
	 \loigiai{
		Thay tọa~độ điểm $M$  vào phương trình đường thẳng $\Delta$, ta có $$-1=\dfrac{3}{m-1}\Leftrightarrow m=-2.$$
	 }
\end{ex}

% nguồn: Phan1_Lop_12_Hinh_hoc_Chuong_5
\begin{ex}%[2H5H2-3]%[ma:603b61e2]
	Trong không gian $Oxyz$, cho các  điểm $A(1;2;3)$, $B(1;1;1)$, $C(3;4;0)$. Đường thẳng $\Delta$ đi qua điểm $A$ và song song với $BC$ có phương trình là
	\choice 
		 {$\dfrac{x+1}{2}=\dfrac{y+2}{3}=\dfrac{z+3}{-1}$}
		 {\True $\dfrac{x-1}{2}=\dfrac{y-2}{3}=\dfrac{z-3}{-1}$}
		 {$\dfrac{x-1}{4}=\dfrac{y-2}{5}=\dfrac{z-3}{1}$}
		{$\dfrac{x+1}{4}=\dfrac{y+2}{5}=\dfrac{z+3}{1}$} 
	 \loigiai{
		Ta có $\overrightarrow{BC}=(2;3;-1)$ nên đường thẳng $\Delta$ có một vectơ chỉ  phương là $\overrightarrow{u}=(2;3;-1)$.\\
		Do đó phương trình của đường thẳng $\Delta$ là $\dfrac{x-1}{2}=\dfrac{y-2}{3}=\dfrac{z-3}{-1}$. 
	} 
\end{ex}

% nguồn: Phan1_Lop_12_Hinh_hoc_Chuong_5
\begin{ex}%[2H5H2-3]%[ma:5718404e]
	Trong không gian $Oxyz$, đường thẳng đi qua  $A(3;2;-1)$ và vuông góc với mặt phẳng $(\alpha)\colon x-2y+z-1=0$ có phương trình chính tắc~là
	\choice 
		{$\dfrac{x-1}{3}=\dfrac{y+2}{2}=\dfrac{z-1}{-1}$} 
		{\True $\dfrac{x-3}{1}=\dfrac{y-2}{-2}=\dfrac{z+1}{1}$} 
		 {$\dfrac{x+1}{3}=\dfrac{y-2}{2}=\dfrac{z+1}{-1}$}
		 {$\dfrac{x+3}{1}=\dfrac{y+2}{-2}=\dfrac{z-1}{1}$}
	\loigiai{ 
		Đường thẳng $\Delta$ đi qua $A(3;2;-1)$ và vuông góc với mặt phẳng $(\alpha)\colon x-2y+z-1=0$ có véc-tơ chỉ  phương~là $\vec{a}=(1;-2;1)$.\\
		Vậy đường thẳng $\Delta$ có phương trình chính  tắc~là $\dfrac{x-3}{1}=\dfrac{y-2}{-2}=\dfrac{z+1}{1}$.
	} 
\end{ex}

% nguồn: Phan1_Lop_12_Hinh_hoc_Chuong_5
\begin{ex}%[2H5H2-3]%[ma:272f0c31]
	Phương trình  nào dưới đây là phương trình chính tắc của đường thẳng\break $d\colon \heva{&x=1+2 t \\ &y=3 t \\ & z=-2+t}$?
	 \choice
		{$\dfrac{x+1}{2}=\dfrac{y}{3}=\dfrac{z-2}{1}$} 
		 {\True $\dfrac{x-1}{2}=\dfrac{y}{3}=\dfrac{z+2}{1}$}
		 {$\dfrac{x+1}{2}=\dfrac{y}{3}=\dfrac{z-2}{-2}$}
		{$\dfrac{x-1}{1}=\dfrac{y}{3}=\dfrac{z+2}{-2}$} 
	\loigiai{ 
		Phương trình chính tắc của đường thẳng $d\colon \heva{&x=1+2 t \\ &y=3 t \\ & z=-2+t}$ là\break $\dfrac{x-1}{2}=\dfrac{y}{3}=\dfrac{z+2}{1}$. 
	 }
\end{ex}

% nguồn: Phan1_Lop_12_Hinh_hoc_Chuong_5
\begin{ex}%[2H5H2-3]%[ma:af52a98f]
	 Trong không gian $Oxyz$, phương trình chính tắc của đường thẳng $AB$ với $A(1; 1; 2)$ và $B(-4; 3;-2)$ là
	 \choice
		 {$\dfrac{x+4}{1}=\dfrac{y-3}{-2}=\dfrac{z+2}{-2}$}
		{\True $\dfrac{x+4}{-5}=\dfrac{y-3}{2}=\dfrac{z+2}{-4}$} 
		 {$\dfrac{x+1}{-5}=\dfrac{y+1}{2}=\dfrac{z+2}{-4}$}
		{$\dfrac{x-1}{1}=\dfrac{y-1}{-2}=\dfrac{z-2}{-2}$} 
	\loigiai{Trong không gian $Oxyz$, phương trình chính tắc của đường thẳng $AB$ với $A(1; 1; 2)$ và  $B(-4; 3;-2)$ là $\dfrac{x+4}{-5}=\dfrac{y-3}{2}=\dfrac{z+2}{-4}$.}
\end{ex}

% nguồn: Phan1_Lop_12_Hinh_hoc_Chuong_5
\begin{ex}%[2H5H2-3]%[ma:e9972608]
	Trong không gian $ Oxyz$, cho  mặt phẳng $(P)\colon x+2y-3z+4=0$ và điểm $ A\left(1;-2;-3\right)$. Phương trình đường thẳng đi qua $ A$ và vuông góc với mặt phẳng $\left({P}\right)$~là
	 \choice
		{\True $\dfrac{x-1}{1}=\dfrac{y+2}{2}=\dfrac{z+3}{-3}$} 
		{$\dfrac{x+1}{1}=\dfrac{y-2}{2}=\dfrac{z-3}{-3}$} 
		 {$\dfrac{x-1}{1}=\dfrac{y-2}{-2}=\dfrac{z+3}{-3}$}
		{$\dfrac{x-1}{1}=\dfrac{y+2}{-2}=\dfrac{z+3}{-3}$} 
	\loigiai{ 
		Gọi $d$~là đường thẳng đi qua $ A$ và vuông góc với mặt phẳng $\left({P}\right)$. Do $d$ vuông góc  với mặt phẳng $\left({P}\right)$ nên nhận véc-tơ $\vec{n}=(1;2;-3)$~làm một véc-tơ chỉ phương.
		Phương trình đường  thẳng $d$~là $\dfrac{x-1}{1}=\dfrac{y+2}{2}=\dfrac{z+3}{-3}$
	} 
\end{ex}

% nguồn: Phan1_Lop_12_Hinh_hoc_Chuong_5
\begin{ex}%[2H5H2-3]%[ma:10e61eb4]
	Trong không gian $ Oxyz$, đường thẳng đi qua  điểm ${A}\left(2;3;1\right)$ và song song với đường thẳng ${d}:\dfrac{{x}-1}{2}=\dfrac{{y}+1}{-4}=\dfrac{{z}-3}{-1}$ có phương trình~là
	\choice 
		 {\True $\dfrac{x-2}{2}=\dfrac{y-3}{-4}=\dfrac{z-1}{-1}$}
		{$ \dfrac{x-2}{2}=\dfrac{y-3}{3}=\dfrac{z-1}{1}$} 
		 {$ \dfrac{x+2}{2}=\dfrac{y+3}{-4}=\dfrac{z+1}{-1}$}
		{$ \dfrac{x-2}{1}=\dfrac{y-3}{-1}=\dfrac{z-1}{3}$} 
	 \loigiai{Gọi $\Delta$~là đường thẳng đi qua điểm $A$ và song song với đường thẳng $d$. Do $\Delta$ song song với $d$ nên đường thẳng $\Delta$ nhận véc-tơ $\vec{u}=(2;-4;-1)$~làm một véc-tơ chỉ phương.
		Phương trình  đường thẳng $\Delta$~là $ \dfrac{x-2}{2}=\dfrac{y-3}{-4}=\dfrac{z-1}{-1}.$
	} 
\end{ex}

% nguồn: Phan1_Lop_12_Hinh_hoc_Chuong_5
\begin{ex}%[2H5H2-3]%[ma:0f7ee65b]
	Trong  không gian $Oxyz$, cho mặt phẳng $(\alpha) \colon -5x+2y+z+2025=0$. Phương trình tham số của đường thẳng đi qua điểm $A(2;1;-5)$ và vuông góc với mặt phẳng $(\alpha)$ là
	 \choice
		 {\True $\heva{&x=2-5t\\&y=1+2t\\&z=-5+t}$}
		{$\heva{&x=5t\\&y=-2t\\&z=-t}$} 
		{$\heva{&x=-5+2t\\&y=2+t\\&z=1-5t}$} 
		{$\heva{&x=-5-2t\\&y=2-t\\&z=1+5t}$} 
	\loigiai{Mặt phẳng $(\alpha)$  có VTPT là $\overrightarrow{n}=(-5;2;1)$. \\
		 Đường thẳng vuông góc với $(\alpha)$ nhận $\overrightarrow{n}=(-5;2;1)$ làm VTCP. \\
		Phương trình tham số của đường  thẳng đi qua điểm $A(2;1;-5)$ và có VTCP $\overrightarrow{n}=(-5;2;1)$ là
		$\heva{&x=2-5t\\&y=1+2t\\&z=-5+t.}$ 
	 }
\end{ex}

% nguồn: Phan1_Lop_12_Hinh_hoc_Chuong_5
\begin{ex}%[2H5H2-3]%[ma:866cfafe]
	Trong không gian $Oxyz$,  cho điểm $A(1;0;3)$. Viết phương trình đường thẳng $\Delta$ đi qua $A$, cắt hai trục $Ox$, $Oz$ tại $M$, $N$ sao cho $\overrightarrow{AM}=2 \overrightarrow{AN}$.
	 \choice
		{$\heva{&x=-1-2t \\ &y=0 \\&  z=3t}$}
		{\True $\heva{&x=-1+2t  \\ &y=0 \\ &z=3t}$}
		 {$\heva{&x=1+t \\ &y=0 \\ &z=3+3t}$}
		{$\heva{&x=1+2t \\  &y=0 \\ &z=3t}$}
	 \loigiai{
		Gọi $M(m;0;0)$, $N(0;0;n)$ là  giao điểm của đường thẳng $\Delta$ với hai trục $Ox$ và $Oz$.\\
		Ta có 
		$\overrightarrow{AM}=(m-1; 0;-3)$; 
		$\overrightarrow{AN}=(0-1; 0-0; n-3)=(-1; 0;  n-3)$.\\
		 Vì $\overrightarrow{AM}=2\overrightarrow{AN}$ nên ta có hệ
		$\heva{&m-1=2 \cdot (-1) \\&0=2 \cdot 0 \\&  -3=2 (n-3)}$
		$\Leftrightarrow \heva{&m=-1  \\&n=\dfrac{3}{2}.}$\\
		Suy ra $M(-1; 0; 0)$ và $N(0; 0;  \dfrac{3}{2})$.\\
		Ta có $\overrightarrow{MN}=(1;  0;\dfrac{3}{2})$.\\
		Véc-tơ chỉ phương của đường  thẳng $\Delta$ là $\overrightarrow{u}=(2;0;3)$.\\
		Vậy phương trình tham số đường thẳng $\Delta$ đi qua điểm $M(-1;0;0)$  có véc-tơ chỉ phương $\overrightarrow{u}=(2;0;3)$ là
		 $\heva{&x=-1+2t\\&y=0\\&z=3t.}$
	} 
\end{ex}

% nguồn: Phan1_Lop_12_Hinh_hoc_Chuong_5
\begin{ex}%[2H5V2-3]%[ma:fc215207]
	Trong không gian  với hệ tọa độ $Oxyz$, phương trình hình chiếu vuông góc của đường thẳng $d\colon\dfrac{x+1}{2}=\dfrac{y-2}{3}=\dfrac{z+3}{1}$ trên mặt phẳng tọa độ $Oxy$ là
	 \choice
		 {$\heva{&x=5-6t\\ &y=11+9t\\ &z=0}$}
		{\True  $\heva{&x=5-6t\\ &y=11-9t\\ &z=0}$}
		{$\heva{&x=5+6t\\ &y=11-9t\\ &z=0}$} 
		 {$\heva{&x=3-6t\\ &y=11-9t\\ &z=0}$}
	\loigiai{ 
		Lấy $N(-1;2;-3)\in d$ và gọi $H$ là hình chiếu của  điểm $N$ trên $(Oxy)$ thì $H(-1;2;0)$.\\
		Phương trình đường thẳng $d$ có dạng  tham số là $\heva{&x=-1+2t\\ &y=2+3t\\ &z=-3+t}$.\\
		Hình  chiếu của $d$ lên $(Oxy)$là đường thẳng $d' \colon \heva{&x=-1+2t\\ &y=2+3t\\ &z=0.}$\\
		Cho $t=3$ ta được điểm  $K(5;11;0)$ thuộc $d'$.\\
		Khi đó  phương trình của $\Delta'$ có thể viết dưới dạng $\heva{&x=5-6t\\ &y=11-9t\\ &z=0.}$\\
	 }
\end{ex}

% nguồn: Phan1_Lop_12_Hinh_hoc_Chuong_5
\begin{ex}%[2H5V2-3]%[ma:9412dbda]
	Trong không  gian với hệ tọa độ $Oxyz$, cho đường thẳng $\Delta\colon\heva{&x=1-t\\ &y=2+t\\ &z=3-2t}$ và mặt phẳng $(P)\colon x-2y+z+5=0$. Phương trình đường thẳng $d$ nằm trong mặt phẳng $(P)$, cắt và vuông góc với $\Delta$ là
	 \choice
		 {$\heva{&x=-3\\ &y=-1+3t\\ &z=1+t}$}
		{$\heva{&x=1-3t\\ &y=2-t\\ &z=-1+t}$} 
		 {$\heva{&x=0\\ &y=-1+5t\\ &z=2+3t}$}
		 {\True $\heva{&x=-3t\\ &y=3-t\\ &z=1+t}$}
	 \loigiai{
		Gọi $I(x;y;z)$ là giao điểm của $\Delta$ và $(P)$. Khi đó tọa độ của  $I$ thỏa mãn $\heva{&x=1-t\\ &y=2+t\\ &z=3-2t\\ &x-2y+z+5=0}$.\\
		Thay $x$, $y$, $z$ từ phương trình $\Delta$ vào phương trình $(P)$, ta được $$(1-t)-2(2+t)+(3-2t)+5=0 \Leftrightarrow 1-t-4-2t+3-2t+5=0 \Leftrightarrow  -5t+5=0 \Leftrightarrow t=1.$$
		Suy  ra $I(0;3;1)$.\\
		 Vectơ chỉ phương của đường thẳng $\Delta$ là $\vec{u}=(-1;1;-2)$.\\
		 Vectơ pháp tuyến của mặt phẳng $(P)$ là $\vec{n}=(1;-2;1)$.\\
		Gọi $\vec{v}$ là vectơ chỉ phương của đường  thẳng $d$.\\
		Vì $d$ nằm trong $(P)$ và  vuông góc với $\Delta$, nên $\vec{v}\perp\vec{u}$ và $\vec{v}\perp\vec{n}$. Do đó $\vec{v}=\left[\vec{u},\vec{n}\right] =(-3;-1;1)$.\\
		Đường thẳng $d$ đi qua điểm  $I(0;3;1)$ và có vectơ chỉ phương $\vec{v}=(-3;-1;1)$ nên có phương trình là $\heva{&x=-3t\\ &y=3-t\\ &z=1+t}$.
	 }
\end{ex}

% nguồn: Phan1_Lop_12_Hinh_hoc_Chuong_5
\begin{ex}%[2H5V2-3]%[ma:fdcaa019]
	Trong  không gian với hệ tọa độ $Oxyz$ cho điểm $A(-1;3;-2)$ và hai mặt phẳng \break{$(P)\colon 2x-y+z+3=0$}; $(Q)\colon x-2y-1=0$. Gọi $d$ là đường thẳng đi qua $A$ và vuông góc với mặt phẳng $(Q)$. Đường thẳng $d'$ là hình chiếu vuông góc của $d$ lên mặt phẳng $(P)$ có phương trình là
	 \choice
		{\True $\heva{&x=u\\ &y=1+4u\\  &z=-2+2u}$}
		{$\heva{&x=u\\  &y=1-7u\\ &z=-2+10u}$}
		 {$\heva{&x=1\\ &y=1+3u\\ &z=-2+u}$}
		 {$\heva{&x=0\\ &y=-1+3u\\ &z=2+u}$}
	\loigiai{ 
		 Mặt phẳng $(Q)$ có vectơ pháp tuyến $\vec{n}_{Q}=(1;-2;0)$.\\
		Vì $d\perp(Q)$ nên $d$ có một vectơ  chỉ phương $\vec{u}_{d}=\vec{n}_{Q}=(1;-2;0)$.\\
		Khi đó phương trình đường thẳng $d$ là  $\heva{&x=-1+t\\ &y=3-2t\\ &z=-2}$.\\
		Gọi $H$ là giao  điểm của đường thẳng $d$ với mặt phẳng $(P)$.\\
		Khi đó, tọa độ điểm $H$ là nghiệm của  hệ phương trình
		$\heva{&x=-1+t\\ &y=3-2t\\ &z=-2\\  &2x-y+z+3=0}$.\\
		Thay  $x$, $y$, $z$ vào phương trình $(P)$, ta được $$2(-1+t)-(3-2t)+(-2)+3=0 \Leftrightarrow 4t-4=0 \Leftrightarrow t=1.$$
		 Suy ra $H(0;1;-2)$.\\
		Để tìm hình chiếu $d'$ của $d$ lên $(P)$, ta tìm mặt phẳng $(R)$ chứa $d$ và vuông góc với $(P)$.\\ 
		Khi đó giao tuyến của $(P)$ và $(R)$ chính  là $d'$.\\
		Vectơ pháp tuyến của $(P)$ là $\vec{n}_P=(2;-1;1)$.\\ 
		Suy ra vectơ pháp tuyến của $(R)$  là $\vec{n}_R=\left[\vec{u}_d,\vec{n}_P\right]=(-2;-1;3)$.\\
		Vì $(R)$ chứa $d$ và $d$ đi  qua $A(-1;3;-2)$ nên phương trình của $(R)$ là $$-2\left(x-(-1)\right)-1(y-3)+3\left(z-(-2)\right)=0 \Leftrightarrow 2x+y-3z-7=0.$$
		Đường thẳng $d'$ là giao tuyến của $(P)\colon2x-y+z+3=0$ và  $(R)\colon 2x+y-3z-7=0$.\\
		 Vectơ chỉ phương của $d'$ là $\left[\vec{n}_P,\vec{n}_R\right]=(2;8;4)$.\\
		Ta chọn vectơ chỉ phương  của $d'$ là $\vec{u'}=\dfrac{1}{2}\left[\vec{n}_P,\vec{n}_R\right]=(1;4;2)$.\\
		Đường thẳng $d'$ đi qua $H$ và có vectơ chỉ phương $\vec{u'}$ có  phương trình là $\heva{&x=u\\ &y=1+4u\\ &z=-2+2u.}$
	} 
\end{ex}

% nguồn: Phan1_Lop_12_Hinh_hoc_Chuong_5
\begin{ex}%[2H5V2-3]%[ma:90f5d714]
	Trong không gian với hệ tọa độ $Oxyz$, cho đường thẳng $\Delta\colon\dfrac{x}{1}=\dfrac{y+1}{2}=\dfrac{z}{1}$ và mặt  phẳng \break{$(P)\colon x+y-z-1=0$}. Đường thẳng nằm trong $(P)$ đồng thời cắt và vuông góc với $\Delta$ có phương trình là
	\choice 
		{\True $\heva{&x=1+3t\\ &y=1-2t\\  &z=1+t}$}
		{$\heva{&x=1+t\\ &y=1-2t\\  &z=1+3t}$}
		 {$\heva{&x=1+3t\\ &y=1-2t\\ &z=2+t}$}
		 {$\heva{&x=1+3t\\ &y=1+2t\\ &z=2+t}$}
	\loigiai{ 
		 Ta có $\Delta\colon\dfrac{x}{1}=\dfrac{y+1}{2}=\dfrac{z}{1}\Rightarrow\Delta\colon\heva{&x=t\\ &y=-1+2t\\ &z=t}$.\\
		Gọi $M=\Delta\cap(P)\Rightarrow M\in\Delta\Rightarrow  M(t;2t-1;t)$.\\
		Mặt  khác $M\in(P)\Rightarrow t+(2t-1)-t-1=0\Leftrightarrow2t-2=0\Leftrightarrow t=1\Rightarrow M(1;1;1)$.\\
		Vectơ pháp tuyến của mặt phẳng $(P)$ là $\vec{n}=(1;1;-1)$.\\ 
		Vectơ chỉ phương của đường  thẳng $\Delta$ là $\vec{u}=(1;2;1)$.\\
		Đường thẳng $d$ nằm trong mặt phẳng $(P)$ đồng thời cắt và vuông góc với $\Delta$ nên đường thẳng $d$ nhận $\left[\vec{n},\vec{u}\right]=(3;-2;1)$ làm vectơ chỉ  phương và $M(1;1;1)\in d$.\\
		Phương trình đường thẳng $d\colon\heva{&x=1+3t\\  &y=1-2t\\ &z=1+t}$.
	 }
\end{ex}

% nguồn: Phan1_Lop_12_Hinh_hoc_Chuong_5
\begin{ex}%[2H5V2-3]%[ma:0857c2c4]
	Trong không gian với hệ toạ độ $Oxyz$, cho hai đường thẳng $d_{1}$ và $d_{2}$ chéo nhau có phương trình $d_{1}\colon\heva{&x=1\\  &y=10+2t\\ &z=t}$ và $d_{2}\colon\heva{&x=3t\\ &y=3-2t\\ &z=-2}$. Gọi $\Delta$ là đường thẳng vuông góc chung của $d_{1}$ và $d_{2}$. Phương trình của $\Delta$ là
	 \choice
		{\True $\heva{&x=2t\\  &y=\frac{177}{98}+3t\\ &z=-6t+\frac{17}{49}}$}{$\heva{&x=\frac{7}{3}+\cdot46t\\ &y=-147t\\ &z=246t}$}
		 {$\heva{&x=1+2t\\ &y=2+3t\\ &z=2-3t}$}
		{$\heva{&x=1+2t\\  &y=2-3t\\ &z=6-4t}$}
	 \loigiai{
		Đường thẳng $d_{1}$ có vectơ chỉ phương là $\vec{u}_{1}=(0;2;1)$, $d_{2}$ có vectơ chỉ phương  là $\vec{u}_{2}=(3;-2;0)$.\\
		Gọi $M(1;10+2t_{1};t_{1})\in d_{1}$, $N(3t_{2};3-2t_{2};-2)\in  d_{2}$.\\
		Suy ra $\vec{MN}=(3t_{2}-1;-2t_{2}-7-2t_{1};-2-t_{1})$.\\ 
		 Ta có
		\begin{eqnarray*}
			\heva{& \vec{MN}\cdot\vec{u}_{1}=0\\  & \vec{MN}\cdot\vec{u}_{2}=0}
			&\Leftrightarrow&\heva{&0(3t_{2}-1)+2(-2t_{2}-7-2t_{1})+1(-2-t_{1})=0\\ &3(3t_{2}-1)-2(-2t_{2}-7-2t_{1})+0(-2-t_{1})=0}\\ 
			 &\Leftrightarrow&\heva{&-5t_{1}-4t_{2}-16=0\\ &4t_{1}+13t_{2}+11=0}\\
			&\Leftrightarrow&\heva{&t_{1}=-\frac{164}{49}\\&t_{2}=\frac{9}{49}.} 
		\end{eqnarray*}
		 Do đó $\vec{MN} = \left(-\dfrac{22}{49}; -\dfrac{33}{49}; \dfrac{66}{49}\right)$.\\
		Ta có  thể chọn một vectơ chỉ phương của $\Delta$ là $\vec{u}=-\dfrac{49}{11}\vec{MN}=(2;3;-6)$.\\
		Từ đó suy ra  phương trình của $\Delta \colon \heva{&x=1+2t\\ &y=\dfrac{162}{49}+3t\\ &z=-\dfrac{164}{49}-6t.}$\\
		Cho $t=-\dfrac{1}{2}$ ta được  điểm $K\left(0;\dfrac{177}{98};\dfrac{17}{49}\right)$ thuộc $\Delta$.\\
		Khi đó phương trình của $\Delta$ có thể viết dưới  dạng $\heva{&x=2t\\ &y=\frac{177}{98}+3t\\ &z=-6t+\frac{17}{49}.}$\\
	 }
\end{ex}

% nguồn: Phan1_Lop_12_Hinh_hoc_Chuong_5
\begin{ex}%[2H5V2-3]%[ma:fca5e35c]
	Trong không gian với hệ tọa độ $Oxyz$, cho đường thẳng $\Delta\colon \dfrac{x}{1}=\dfrac{y+1}{2}=\dfrac{z-1}{1}$ và mặt phẳng $(P)\colon x-2y-z+3=0$. Đường thẳng nằm trong $(P)$ đồng thời cắt và vuông  góc với $\Delta$ có phương trình là
	 \choice
		{$\heva{&x=1+2t\\  &y=1-t\\ &z=2}$}
		{$\heva{&x=-3\\  &y=-t\\ &z=2t}$}
		{$\heva{&x=1+t\\  &y=1-2t\\ &z=2+3t}$}
		{\True $\heva{&x=1\\  &y=1-t\\ &z=2+2t}$}
	\loigiai{ 
		Ta có  $\Delta\colon\dfrac{x}{1}=\dfrac{y+1}{2}=\dfrac{z-1}{1}\Rightarrow\Delta\colon\heva{&x=t\\ &y=-1+2t\\ &z=1+t}$.\\
		Gọi $M=\Delta\cap(P)\Rightarrow M\in\Delta\Rightarrow  M(t;2t-1;t+1)$.\\
		Mặt khác $M\in(P)\Rightarrow  t-2(2t-1)-(t+1)+3=0\Leftrightarrow t-4t+2-t-1+3=0\\ \Leftrightarrow-4t+4=0\Leftrightarrow t=1\Rightarrow M(1;1;2)$.\\
		Vectơ pháp tuyến của  mặt phẳng $(P)$ là $\vec{n}=(1;-2;-1)$.\\
		Vectơ  chỉ phương của đường thẳng $\Delta$ là $\vec{u}=(1;2;1)$.\\
		Đường thẳng $d$ nằm trong mặt phẳng $(P)$ đồng thời cắt và vuông góc với $\Delta$.\\ 
		Và $[\vec{n},\vec{u}]=(0;-2;4)$ nên đường thẳng $d$ nhận $\vec{u}_d=\dfrac{1}{2}\left[\vec{n},\vec{u}\right]=(0;-1;2)$ làm vectơ chỉ  phương.\\
		Mà $M(1;1;2)\in d$ nên phương trình  đường thẳng $d\colon\heva{&x=1\\ &y=1-t\\ &z=2+2t}$.
	} 
\end{ex}

% nguồn: Phan1_Lop_12_Hinh_hoc_Chuong_5
\begin{ex}%[2H5V2-3]%[ma:5d5ddf42]
	Trong không gian với hệ tọa độ $Oxyz$,  cho điểm $M(0;2;0)$ và đường thẳng \break{$d\colon\heva{&x=4+3t\\ &y=2+t\\ &z=-1+t},t\in\mathbb{R}$}.
	Đường thẳng đi qua $M$, cắt  và vuông góc với $d$ có phương trình là
	\choice 
		 {$\dfrac{x}{-1}=\dfrac{y}{1}=\dfrac{z-1}{2}$}
		{\True $\dfrac{x}{-1}=\dfrac{y-2}{1}=\dfrac{z}{2}$}{$\dfrac{x-1}{1}=\dfrac{y}{-1}=\dfrac{z}{-2}$}{$\dfrac{x-1}{1}=\dfrac{y-1}{1}=\dfrac{z}{2}$} 
	\loigiai{ 
		Ta có đường thẳng  $d$ đi qua điểm $N(4;2;-1)$ và có vectơ chỉ phương $\vec{u}_{d}=(3;1;1)$.\\
		Gọi $H$ là hình chiếu vuông góc của $M$ lên đường thẳng  $d$.\\
		Suy ra $H\in d$  hay $H(4+3t;2+t;-1+t)$, suy ra $\vec{MH}=(4+3t;t;-1+t)$.\\
		Ta có 
		$\vec{MH}\cdot\vec{u}_{d}=0\Leftrightarrow3(4+3t)+1(t)+1(-1+t)=0\Leftrightarrow11t+11=0\Leftrightarrow t=-1\\ \Rightarrow H(1;1;-2)$.\\ 
		Đường thẳng đi qua $M(0;2;0)$ và vuông góc với $d$ có vectơ chỉ  phương là $\vec{HM}=(-1;1;2)$ có phương trình là $\dfrac{x}{-1}=\dfrac{y-2}{1}=\dfrac{z}{2}$.
	 }
\end{ex}

% nguồn: Phan1_Lop_12_Hinh_hoc_Chuong_5
\begin{ex}%[2H5V2-3]%[ma:c2554514]
	Trong không gian với hệ toạ độ $Oxyz$, cho đường thẳng  $d\colon\dfrac{x+2}{1}=\dfrac{y-1}{1}=\dfrac{z-2}{2}$. Đường thẳng $d'$ là hình chiếu của $d$ lên mặt phẳng $Oxy$, có phương trình là
	 \choice
		{\True $\heva{&x=-3+t\\ &y=t\\  &z=0}$}
		{$\heva{&x=-3+t\\  &y=-t\\ &z=0}$}
		{$\heva{&x=-3+t\\ &y=1+t\\  &z=0}$}
		 {$\heva{&x=3-t\\ &y=-t\\ &z=0}$}
	 \loigiai{
		Phương trình tham số của $d\colon\heva{&x=-2+t\\ &y=1+t\\ &z=2+2t}$.\\ 
		Hình chiếu của $d$  lên mặt phẳng $(Oxy)$ là đường thẳng $\Delta' \colon \heva{&x=-2+t\\ &y=1+t\\ &z=0}$.\\
		Cho $t=-1$  ta được điểm $B(-3;0;0)$ thuộc $\Delta'$.\\
		Khi đó  phương trình của $\Delta' \colon \heva{&x=-3+t\\ &y=t\\ &z=0}$.\\
	} 
\end{ex}

% nguồn: Phan1_Lop_12_Hinh_hoc_Chuong_5
\begin{ex}%[2H5V2-3]%[ma:17a56e90]
	Trong không  gian với hệ toạ độ $Oxyz$, cho hai điểm $A(3;3;1)$, $B(3;2;1)$, mặt phẳng $(P)\colon x-2y-z+3=0$. Đường thẳng $d$ nằm trên $(P)$ sao cho mọi điểm của $d$ cách đều hai điểm $A$, $B$ có phương trình là
	\choice 
		{\True  $\heva{&x=4+t\\ &y=\frac{5}{2}\\ &z=2+t}$}
		{$\heva{&x=t\\ &y=7+3t\\ &z=2t}$} 
		{$\heva{&x=4+t\\  &y=\frac{5}{2}+2t\\ &z=2+t}$}
		{$\heva{&x=4+t\\  &y=-\frac{5}{2}\\ &z=2+t}$}
	\loigiai{ 
		Ta có $\vec{AB}=(0;-1;0)$ và $I\left(3;\dfrac{5}{2};1\right)$ là trung  điểm của $AB$.\\
		Gọi $(\alpha)$ là mặt phẳng trung  trực của $AB$. Khi đó $d=(\alpha)\cap(P)$.\\
		Mặt phẳng $(\alpha)$ đi qua $I$ và có vectơ pháp tuyến $\vec{AB}=(0;-1;0)$ có phương trình là $$0(x-3)+(-1)\left(y-\frac{5}{2}\right)+0(z-1)=0\Leftrightarrow -y+\frac{5}{2}=0\Leftrightarrow  2y-5=0.$$
		Vì $d$ là đường giao tuyến của $(\alpha)$ và $(P)$ nên một vectơ chỉ phương của  $d$ là $$\vec{u}_{d}=\left[\vec{n}_{(P)},\vec{n}_{(\alpha)}\right]=(2;0;2).$$
		Ta có thể chọn $\vec{u}=\dfrac{1}{2}\vec{u}_{d}=(1;0;1)$ làm vectơ chỉ phương  của $d$.\\
		Để tìm một điểm thuộc $d$, ta giải hệ phương trình của  hai mặt phẳng
		 $$\heva{&x-2y-z+3=0\\ &2y-5=0.}$$
		 Từ phương trình thứ hai, suy ra $y=\dfrac{5}{2}$. Thay vào phương trình thứ nhất, ta được $$x-5-z+3=0 \Leftrightarrow x-z-2=0.$$
		Chọn $x=4$,  thì $z=x-2=4-2=2$, suy ra điểm $C\left(4;\dfrac{5}{2};2\right)$ thuộc $d$.\\
		Vậy $d$ có phương trình tham số là $\heva{&x=4+t\\ &y=\frac{5}{2}\\ &z=2+t}$. 
	} 
\end{ex}

% nguồn: Phan1_Lop_12_Hinh_hoc_Chuong_5
\begin{ex}%[2H5V2-3]%[ma:defedf5b]
	 Trong không gian với hệ tọa độ $Oxyz$, gọi $(\alpha)$ là mặt phẳng chứa đường thẳng $\Delta$ có phương trình $\dfrac{x-2}{1}=\dfrac{y-1}{1}=\dfrac{z}{2}$ và vuông góc với mặt phẳng $(\beta)\colon x+y-2z-1=0$. Giao tuyến của $(\alpha)$ và $(\beta)$ đi
	qua điểm nào trong  các điểm sau?
	 \choice
		 {$C(1;2;1)$}
		 {$D(2;1;0)$}
		{$B(0;1;0)$} 
		{\True  $A(2;1;1)$}
	 \loigiai{
		Ta có vectơ chỉ phương của đường thẳng $\Delta$ là  $\vec{u}=(1;1;2)$.\\
		Vectơ  pháp tuyến của mặt phẳng $(\beta)$ là $\vec{n}_{\beta}=(1;1;-2)$.\\
		Vì $(\alpha)$ là mặt phẳng chứa  đường thẳng $\Delta$ và vuông góc với mặt phẳng $(\beta)$, nên $(\alpha)$ có một vectơ pháp tuyến là
		 $\vec{n}_{\alpha}=\left[\vec{u},\vec{n}_{\beta}\right]=(-4;4;0)$.\\
		Ta có thể chọn $\vec{a}=(1;-1;0)$ là vectơ pháp tuyến  của $(\alpha)$.\\
		 Đường thẳng $\Delta$ đi qua điểm $M(2;1;0)$.\\
		Phương trình mặt phẳng $(\alpha)$  là $1(x-2)-1(y-1)+0(z-0)=0 \Leftrightarrow x-y-1=0$.\\
		Gọi  $d=(\alpha)\cap(\beta)$.\\
		 Vectơ chỉ phương của $d$ là $\left[\vec{n}_{\alpha},\vec{n}_{\beta}\right]=(2;2;2)$.\\
		Ta có thể chọn $\vec{u}_{d}=(1;1;1)$ là vectơ chỉ  phương của $d$.\\
		Giao điểm của  đường thẳng $\Delta$ và mặt phẳng $(\beta)$ là $I(3;2;2)$.\\
		Vì $\heva{&I(3;2;2) \in \Delta\\&\Delta \subset (\alpha)\\&I \in (\beta)}$ nên $I$ là một điểm  chung của $(\alpha)$ và $(\beta)$, tức là $I \in d$.\\
		Suy ra phương trình đường thẳng $d\colon\dfrac{x-3}{1}=\dfrac{y-2}{1}=\dfrac{z-2}{1}$.\\ 
		Thay toạ độ điểm $A(2;1;1)$ vào phương trình đường thẳng $d$ ta được $\dfrac{2-3}{1}=\dfrac{1-2}{1}=\dfrac{1-2}{1}$ (đúng).\\ 
		Vậy $A(2;1;1)$  thuộc đường thẳng $d$.
	 }
\end{ex}

% nguồn: Phan1_Lop_12_Hinh_hoc_Chuong_5
\begin{ex}%[2H5V2-3]%[ma:c369f689]
	Trong không gian với hệ tọa độ $Oxyz$, đường vuông góc chung của hai đường thẳng $d\colon\heva{&x=1+t\\ &y=0\\ &z=-5+t}$ và  $d'\colon\heva{&x=0\\ &y=4-2t'\\ &z=5+3t'}$ có phương trình là
	\choice 
		{$\dfrac{x+4}{-2}=\dfrac{y}{3}=\dfrac{z-2}{2}$} 
		{\True $\dfrac{x-4}{-2}=\dfrac{y}{3}=\dfrac{z+2}{2}$}{$\dfrac{x-4}{-1}=\dfrac{y}{3}=\dfrac{z+2}{1}$}{$\dfrac{x-4}{2}=\dfrac{y}{-3}=\dfrac{z-2}{-2}$} 
	 \loigiai{
		Giả sử $AB$ là đường vuông góc chung của $d$ và $d'$ với $A\in d$,  $B\in d'$.\\
		Ta có $\vec{u}_{d}=(1;0;1)$, $\vec{u}_{d'}=(0;-2;3)$ lần lượt là vectơ chỉ  phương của $d$ và $d'$.\\
		Ta có $A(1+a;0;-5+a)$ và $B(0;4-2b;5+3b)$. Suy  ra $\vec{BA}=(1+a;2b-4;a-3b-10)$.\\
		Khi đó 
		\begin{eqnarray*}
			\heva{&\vec{u}_{d}\cdot\vec{BA} =  0\\ & \vec{u}_{d'} \cdot \vec{BA}=0}
			 &\Leftrightarrow&\heva{&1(1+a)+0(2b-4)+1(a-3b-10)=0\\ &0(1+a)-2(2b-4)+3(a-3b-10)=0}\\
			 &\Leftrightarrow&\heva{&2a-3b-9=0\\ &3a-13b-22=0}\\
			 &\Leftrightarrow&\heva{&a=3\\ &b=-1.}
		\end{eqnarray*}
		Suy ra  $A(4;0;-2)$, $B(0;6;2)$ và $\vec{AB}=(-4;6;4)$.\\
		Ta có thể chọn $\vec{u}=\dfrac{1}{2}\vec{AB}=(-2;3;2)$  là một vectơ chỉ phương của $AB$.\\
		Vậy đường thẳng $AB$ đi qua điểm  $A(4;0;-2)$ và có vectơ chỉ phương $\vec{u}=(-2;3;2)$ nên có phương trình chính tắc là $\dfrac{x-4}{-2}=\dfrac{y}{3}=\dfrac{z+2}{2}$.
	 }
\end{ex}

% nguồn: Phan1_Lop_12_Hinh_hoc_Chuong_5
\begin{ex}%[2H5V2-3]%[ma:83934f77]
	Trong không gian với hệ tọa độ $Oxyz$, cho hai  đường thẳng chéo nhau \break{$d_{1}\colon\dfrac{x-2}{1}=\dfrac{y-1}{-1}=\dfrac{z-2}{-1}$} và $d_{2}\colon\heva{&x=t\\ &y=3\\ &z=-2+t}$. Phương trình đường vuông góc chung của hai đường thẳng $d_{1}$ và $d_{2}$ là
	 \choice
		{\True $\heva{&x=2+t\\ &y=1+2t\\  &z=2-t}$}
		 {$\heva{&x=3+t\\ &y=3-2t\\ &z=t-t}$}
		{$\heva{&x=2+3t\\  &y=1-2t\\ &z=2-5t}$}
		{$\heva{&x=3+t\\ &y=3\\  &z=t-t}$}
	 \loigiai{
		Gọi  $d$ là đường thẳng cần tìm.\\
		Gọi $A=d\cap d_{1}$, $B=d\cap d_{2}$.\\ 
		Khi đó $A\in d_{1}\Rightarrow A(2+a;1-a;2-a)$; $B\in d_{2}\Rightarrow B(b;3;-2+b)$  và $\vec{AB}=(-a+b-2;a+2;a+b-4)$.\\
		Mà $d_{1}$ có vectơ  chỉ phương $\vec{u}_{1}=(1;-1;-1)$, $d_{2}$ có vectơ chỉ phương $\vec{u}_{2}=(1;0;1)$.\\
		Khi  đó
		\begin{eqnarray*}
			\heva{&\vec{AB}\cdot\vec{u}_{1}=0\\ &  \vec{AB}\cdot\vec{u}_{2}=0}
			&\Leftrightarrow&\heva{&1(-a+b-2)-1(a+2)-1(a+b-4)=0\\ &1(-a+b-2)+0(a+2)+1(a+b-4)=0}\\ 
			&\Leftrightarrow&\heva{&-3a=0\\ &2b-6=0}\\ 
			&\Leftrightarrow&\heva{&a=0\\  &b=3.}
		\end{eqnarray*}
		Suy ra $A(2;1;2)$ và  $B(3;3;1)$.\\
		Vậy $d$ đi qua $A$ và có vectơ  chỉ phương $\vec{AB}=(1;2;-1)$ nên phương trình của $d \colon \heva{&x=2+t\\ &y=1+2t\\ &z=2-t}$.
	} 
\end{ex}

% nguồn: Phan1_Lop_12_Hinh_hoc_Chuong_5
\begin{ex}%[2H5V2-3]%[ma:b5c048a8]
	Trong không gian với hệ tọa độ $Oxyz$  cho điểm $A(1;0;2)$ và đường thẳng $d$ có phương trình $\dfrac{x-1}{1}=\dfrac{y}{1}=\dfrac{z+1}{2}$. Phương trình đường thẳng $\Delta$ đi qua $A$, vuông góc và cắt $d$ là
	\choice 
		{$\dfrac{x-1}{2}=\dfrac{y}{2}=\dfrac{z-2}{1}$}{$\dfrac{x-1}{1}=\dfrac{y}{-3}=\dfrac{z-2}{1}$}{$\dfrac{x-1}{1}=\dfrac{y}{1}=\dfrac{z-2}{1}$} 
		{\True $\dfrac{x-1}{1}=\dfrac{y}{1}=\dfrac{z-2}{-1}$} 
	 \loigiai{
		Đường thẳng $d$ đi  qua điểm $M(1;0;-1)$ và có vectơ chỉ phương $\vec{u}=(1;1;2)$ nên có phương trình tham số là $\Delta\colon\heva{&x=1+t\\ &y=t\\ &z=-1+2t}$.\\
		Gọi $(P)$ là  mặt phẳng qua điểm $A$ và vuông góc với đường thẳng $d$, nên nhận vectơ chỉ phương của $d$ là vectơ pháp tuyến, suy ra $(P)\colon 1(x-1)+1(y-0)+2(z-2)=0\Leftrightarrow x+y+2z-5=0$.\\
		Gọi $B$  là giao điểm của mặt phẳng $(P)$ và đường thẳng $d$.\\
		 Suy ra $B\in d$ hay $B(1+t;t;-1+2t)$.\\
		Mà $B\in (P)\Leftrightarrow(1+t)+t+2(-1+2t)-5=0\Leftrightarrow 1+t+t-2+4t-5=0\Leftrightarrow 6t-6=0\Leftrightarrow t=1$  hay $B(2;1;1)$.\\
		Đường thẳng $\Delta$ đi qua $A(1;0;2)$ và nhận  vectơ $\vec{AB}=(1;1;-1)$ làm vectơ chỉ phương có phương trình là $\Delta\colon\dfrac{x-1}{1}=\dfrac{y}{1}=\dfrac{z-2}{-1}$.
	} 
\end{ex}

% nguồn: Phan1_Lop_12_Hinh_hoc_Chuong_5
\begin{ex}%[2H5V2-3]%[ma:a5291d1c]
	Trong không gian với hệ tọa độ $Oxyz$, cho điểm $A(0;1;-1)$ và đường thẳng \break{$d\colon\dfrac{x+3}{4}=\dfrac{y-1}{-1}=\dfrac{z-3}{-4}$}. Đường thẳng $\Delta$ đi  qua điểm $A$, vuông góc và cắt đường thẳng $d$ có phương trình là
	 \choice
		{$\dfrac{x}{-13}=\dfrac{y-1}{28}=\dfrac{z+1}{20}$} 
		{\True $\dfrac{x}{13}=\dfrac{y-1}{-28}=\dfrac{z+1}{20}$}{$\dfrac{x}{13}=\dfrac{y-1}{28}=\dfrac{z+1}{-20}$}{$\dfrac{x}{13}=\dfrac{y-1}{28}=\dfrac{z+1}{20}$} 
	 \loigiai{
		Gọi $B$ là giao điểm của đường thẳng $d$ và  đường thẳng $\Delta$.\\
		Đường thẳng  $d$ có phương trình tham số $\heva{&x=-3+4t\\ &y=1-t,\\ &z=3-4t}$, $t\in\mathbb{R}$.\\
		 Và $B\in d\Rightarrow B(-3+4t;1-t;3-4t)$, suy ra $\vec{AB}=(-3+4t;-t;4-4t)$.\\
		Đường thẳng  $d$ có vectơ chỉ phương $\vec{u}=(4;-1;-4)$.\\
		Ta có  $\vec{AB}\cdot\vec{u}=0\Leftrightarrow4(-3+4t)-1(-t)-4(4-4t)=0\Leftrightarrow33t=28\Leftrightarrow t=\dfrac{28}{33}$.\\
		Suy  ra $\vec{AB}=\left(\dfrac{13}{33};\dfrac{-28}{33};\dfrac{20}{33}\right)$.\\
		Ta chọn vectơ  chỉ phương của $\Delta$ là $\vec{u}_{\Delta}=33\vec{AB}=(13;-28;20)$.\\
		Đường thẳng $\Delta$ đi qua điểm $A(0;1;-1)$ và nhận vectơ  $\vec{AB}$ hay $\vec{u}_{\Delta}=(13;-28;20)$ là vectơ chỉ phương, có phương trình chính tắc là $\dfrac{x}{13}=\dfrac{y-1}{-28}=\dfrac{z+1}{20}$.
	} 
\end{ex}

% nguồn: Phan1_Lop_12_Hinh_hoc_Chuong_5
\begin{ex}%[2H5V2-3]%[ma:07902a3d]
	Trong không gian với hệ tọa độ $Oxyz$, cho đường thẳng $\Delta\colon\heva{&x=-5+t\\ &y=2t\\ &z=1+7t}$ và mặt phẳng $(P)\colon  3x+y-2z-1=0$. Phương trình đường thẳng $d$ nằm trong mặt phẳng $(P)$, cắt và vuông góc với $\Delta$ là
	 \choice
		{\True $\heva{&x=-7-11t\\  &y=-4+23t\\ &z=-13-5t}$}
		{$\heva{&x=-11-7t\\ &y=23-4t\\ &z=-5-13t}$} 
		 {$\heva{&x=3-t\\ &y=-1+5t\\&z=2+3t}$}
		{$\heva{&x=4-11t\\ &y=3+23t\\  &z=1-5t}$}
	\loigiai{ 
		Gọi $I(x;y;z)$ là giao điểm của $\Delta$  và $(P)$. Khi đó tọa độ của $I$ thỏa mãn
		$\heva{&x=-5+t\\ &y=2t\\ &z=1+7t\\ &3x+y-2z-8=0.}$\\ 
		Thay $x$, $y$, $z$ từ phương trình $\Delta$ vào phương trình $(P)$,  ta được $$3(-5+t)+(2t)-2(1+7t)-1=0 \Leftrightarrow -15+3t+2t-2-14t-1=0 \Leftrightarrow -9t-18=0 \Leftrightarrow t=-2.$$
		 Suy ra $I(-7;-4;-13)$.\\
		Vectơ chỉ phương của đường thẳng $\Delta$ là  $\vec{u}=(1;2;7)$.\\
		 Vectơ pháp tuyến của mặt phẳng $(P)$ là $\vec{n}=(3;1;-2)$.\\
		Gọi $\vec{v}$ là vectơ  chỉ phương của đường thẳng $d$.\\
		Vì $d$ nằm trong $(P)$ và vuông  góc với $\Delta$, nên $\vec{v} \perp \vec{u}$ và $\vec{v} \perp \vec{n}$.\\
		 Do đó $\vec{v} = \left[\vec{u},\vec{n}\right] = (-11; 23; -5)$.\\
		Đường thẳng $d$ đi qua điểm $I(-7;-4;-13)$ và có vectơ chỉ phương $\vec{v}=(-11;23;-5)$ nên có  phương trình là $\heva{&x=-7-11t\\ &y=-4+23t\\ &z=-13-5t}$.
	} 
\end{ex}

% nguồn: Phan1_Lop_12_Hinh_hoc_Chuong_5
\begin{ex}%[2H5V2-3]%[ma:1d46e84c]
	Trong không gian với hệ  tọa độ $Oxyz$, cho điểm $M(2;1;0)$ và đường thẳng \break{$\Delta\colon\dfrac{x-1}{2}=\dfrac{y+1}{1}=\dfrac{z}{-1}$}. Phương trình của đường thẳng $d$ đi qua điểm $M$, cắt và vuông góc với $\Delta$ là
	\choice 
		 {$\dfrac{x-2}{1}=\dfrac{y-1}{-4}=\dfrac{z}{1}$}{$\dfrac{x-2}{1}=\dfrac{y-1}{4}=\dfrac{z}{1}$}{$\dfrac{x-2}{2}=\dfrac{y-1}{-4}=\dfrac{z}{1}$}
		{\True  $\dfrac{x-2}{1}=\dfrac{y-1}{-4}=\dfrac{z}{-2}$}
	 \loigiai{
		Gọi  $N=d\cap\Delta\Rightarrow N\in\Delta$ nên $N(1+2t;-1+t;-t)$.\\
		Khi đó ta có $\vec{MN}=(2t-1;t-2;-t).$\\ 
		Đường thẳng $\Delta$ có  vectơ chỉ phương $\vec{a}=(2;1;-1)$.\\
		Vì $\vec{MN}\cdot\vec{a}=0\Leftrightarrow2(2t-1)+1(t-2)+(-1)(-t)=0\Leftrightarrow  6t-4=0\Leftrightarrow t=\dfrac{2}{3}\Rightarrow\vec{MN}=\left(\dfrac{1}{3};-\dfrac{4}{3};-\dfrac{2}{3}\right)$.\\
		 Ta chọn vectơ chỉ phương của $d$ là $\vec{a}_{d}=3\vec{MN}=(1;-4;-2)$.\\
		Vậy phương trình của  đường thẳng $d\colon\dfrac{x-2}{1}=\dfrac{y-1}{-4}=\dfrac{z}{-2}$.
	 }
\end{ex}

% nguồn: Phan1_Lop_12_Hinh_hoc_Chuong_5
\begin{ex}%[2H5V2-3]%[ma:72af3ba1]
	Trong không gian với hệ toạ độ $Oxyz$, gọi $(\Delta)$ là đường vuông góc  chung của hai đường thẳng $(d_{1})\colon \dfrac{x-1}{-1}=\dfrac{y-2}{2}=\dfrac{z}{3}$ và $(d_{2})\colon\heva{&x=1+t'\\ &y=3-2t'\\ &z=1}$. Phương trình của $\Delta$ là
	\choice 
		 {\True $\heva{&x=\frac{2}{3}+21\\ &y=\frac{8}{3}+1\\ &z=1}$}
		{$\heva{&x=\frac{2}{3}=2t\\ &y=\frac{8}{3}+t\\  &z=1}$}
		{$\heva{&x=\frac{2}{3}+2t\\ &y=\frac{88}{3}+t\\  &z=1-2t}$}{$\heva{&x=\frac{2}{3}+2t\\ &y=\frac{8}{3}+t\\ &z=1+2t}$}
	\loigiai{ 
		Đường thẳng $d_1$ có vectơ chỉ phương là  $\vec{u}_1=(-1;2;3)$, $d_2$ có vectơ chỉ phương là $\vec{u}_2=(1;-2;0)$.\\
		Gọi $M$, $N$ lần lượt là giao điểm của $\Delta$ với hai đường thẳng  $d_{1}$, $d_{2}$.\\
		Suy ra $M(1-t;2+2t;3t)\in  d_{1}$ và $N(1+t';3-2t';1)\in d_{2}$.\\
		Khi đó  $\vec{MN}=(t'+t;1-2t'-2t;1-3t)$.\\
		Ta có 
		\begin{eqnarray*}
			\heva{& \vec{MN}\cdot\vec{u}_{1}=0\\ &  \vec{MN}\cdot\vec{u}_{2}=0}
			 &\Leftrightarrow&\heva{&(-1)(t'+t)+2(1-2t'-2t)+3(1-3t)=0\\ &1(t'+t)-2(1-2t'-2t)+0(1-3t)=0}\\
			 &\Leftrightarrow&\heva{&-5t'-14t+5=0\\ &5t'+5t-2=0}\\
			&\Leftrightarrow&\heva{&t=\frac{1}{3}\\&t'=\frac{1}{15}.} 
		\end{eqnarray*}
		Do đó $\vec{MN}=\left(\dfrac{2}{5};\dfrac{1}{5};0\right)$ và  $M\left(\dfrac{2}{3};\dfrac{8}{3};1\right)$.\\
		Ta có thể chọn một  vectơ chỉ phương của $\Delta$ là $\vec{u}=5\vec{MN}=(2;1;0)$.\\
		Từ  đó suy ra phương trình của $\Delta\colon\heva{&x=\frac{2}{3}+2t\\ &y=\frac{8}{3}+t\\ &z=1}$.
	} 
\end{ex}

% nguồn: Phan1_Lop_12_Hinh_hoc_Chuong_5
\begin{ex}%[2H5V2-3]%[ma:e53f8c2f]
	Trong không gian với hệ tọa độ $Oxyz$, cho ba mặt phẳng $(Oxy)$; $(Oyz)$; $(P)\colon x+2y-4=0$. Gọi $d$ là giao tuyến của hai mặt phẳng $(Oxy)$ và $(Oyz)$. Đường thẳng $d'$  là hình chiếu vuông góc của $d$ lên mặt phẳng $(P)\colon x+2y-4=0$ có phương trình là
	 \choice
		{\True $\heva{&x=-2u\\  &y=2+u\\ &z=0}$}
		{$\heva{&x=2u\\ &y=2+u\\  &z=0}$}
		{$\heva{&x=-2u\\ &y=2+u\\ &z=u}$} 
		{$\heva{&x=-2u\\  &y=u\\ &z=2}$}
	 \loigiai{
		Vì  $d$ là giao tuyến của hai mặt phẳng $(Oxy)$ và $(Oyz)$ nên $d$ chính là trục $Oy\colon\heva{&x=0\\ &y=t\\ &z=0.}$\\
		Gọi $H$ là giao điểm của đường  thẳng $d$ với mặt phẳng $(P)$, toạ độ điểm $H$ là nghiệm của hệ phương trình
		$\heva{&x=0\\ &y=t\\ &z=0\\  &x+2y-4=0.}$ \\
		Thay $x$, $y$, $z$  vào phương trình $(P)$, ta được $0+2t-4=0 \Leftrightarrow 2t=4 \Leftrightarrow t=2$.\\
		Suy ra $H(0;2;0)$.\\ 
		Để tìm hình chiếu $d'$ của $d$ lên $(P)$, ta tìm  mặt phẳng $(R)$ chứa $d$ và vuông góc với $(P)$.\\
		Khi đó giao tuyến của  $(P)$ và $(R)$ chính là $d'$.\\
		 Vectơ pháp tuyến của $(P)$ là $\vec{n}_P=(1;2;0)$.\\
		Vectơ pháp tuyến của $(R)$ là  $\vec{n}_R = \left[\vec{u}_d,\vec{n}_P\right]=(0;0;-1)$.\\
		Vì $(R)$ chứa $d$ (trục $Oy$) nên $(R)$ đi  qua $O(0;0;0)$ hay phương trình của $(R)$ là
		$$0(x-0)+0(y-0)+1(z-0)=0 \Leftrightarrow  z=0.$$
		Suy ra $d'$ là giao tuyến của $(P)\colon  x+2y-4=0$ và $(R)\colon z=0$.\\
		Vectơ chỉ phương của $d'$ là  $\left[\vec{n}_P, \vec{n}_R\right]=(2;-1;0)$.\\
		Ta chọn $\vec{u'}=-\left[\vec{n}_P, \vec{n}_R\right]=(-2;1;0)$ làm vectơ chỉ phương của $d'$.\\ 
		Đường thẳng $d'$ đi qua $H$ và có vectơ chỉ phương $\vec{u'}$ có  phương trình là $\heva{&x=-2u\\ &y=2+u\\ &z=0},\ u \in\mathbb{R}$.
	 }
\end{ex}

% nguồn: Phan1_Lop_12_Hinh_hoc_Chuong_5
\begin{ex}%[2H5V2-3]%[ma:9444a6f7]
	Trong  không gian với hệ trục tọa độ $Oxyz$, cho điểm $A(1;3;2)$, $B(3;3;0)$ và đường thẳng $d\colon\dfrac{x}{-2}=\dfrac{y+2}{1}=\dfrac{z-1}{3}$. Gọi $M$ là trung điểm của $AB$. Đường thẳng $\Delta$ đi qua $H(1;1;1)$ và cắt hai đường thẳng $d$ và $OM$ có phương trình là
	\choice 
		 {$\heva{&x=1+2t\\ &y=1+23t\\ &z=1-15t}$}
		{$\heva{&x=1+2t\\  &y=1+23t\\ &z=1+15t}$}
		{$\heva{&x=1-2t\\ &y=1-23t\\ &z=1-15t}$} 
		{\True $\heva{&x=1+2t\\  &y=1-23t\\ &z=1-15t}$}
	 \loigiai{
		 Tọa độ trung điểm của đoạn thẳng $AB$ là $M(2;3;1)$.\\
		Gọi $(P)$ là mặt  phẳng đi qua $H$ và chứa $d$.\\
		 Đường thẳng $d$ đi qua $C(0;-2;1)$ và có vectơ chỉ phương $\vec{u}=(-2;1;3)$.\\
		 Và $\vec{HC}=(-1;-3;0)$.\\
		Mặt phẳng $(P)$ có vectơ pháp tuyến là  $\vec{n}_{p}=\left[\vec{u},\vec{HC}\right]=(9;3;7)$.\\
		Gọi $(Q)$ là mặt phẳng đi qua $H$  và chứa $OM$.\\
		Đường thẳng  $OM$ đi qua $O(0;0;0)$ và có vectơ chỉ phương $\vec{OM}=(2;3;1)$. \\
		Và  $\vec{OH}=(1;1;1)$.\\
		Mặt phẳng $(Q)$ có  vectơ pháp tuyến là $\vec{n}_{Q}=\left[\vec{OM},\vec{OH}\right]=(2;-1;-1)$.\\
		Đường thẳng  $\Delta$ đi qua $H(1;1;1)$ và có vectơ chỉ phương $\left[\vec{n}_{p},\vec{n}_{Q}\right]=(2;23;-15)$ nên có phương trình
		$\heva{&x=1+2t\\ &y=1+23t\\ &z=1-15t}$. 
	 }
\end{ex}


%%%%%%%%%%%%%%%% ĐÚNG/SAI %%%%%%%%%%%%%%%%

% nguồn: Phan1_Lop_12_Hinh_hoc_Chuong_5
\begin{ex}%[2H5H2-3]%[ma:42261c5c]
	Trong không gian với hệ tọa độ $Oxyz$, cho hai điểm $A(-3;1;2)$, $B(1;5;-2)$ và mặt phẳng $(P)$ có phương trình tổng quát  là $2x - 4y + 4z + 3 = 0$.
	\choiceTF 
		{Khoảng cách từ  điểm $A$ đến mặt phẳng $(P)$ bằng $\dfrac{1}{6}$}
		{Gọi đường thẳng $d$ đi qua $A$ và vuông góc với $(P)$. Phương trình tham số của đường thẳng $d$ là $\heva{&x = -3 +  t \\ &y = 1 - 2t \\ &z = 2 + 2t}, (t \in \mathbb{R})$}
		{Gọi $(\alpha)$ là mặt phẳng đi qua hai điểm $A$, $B$ đồng thời  vuông góc với mặt phẳng $(P)$. Mặt phẳng $(\alpha)$ có phương trình tổng quát là $5x - 4y + z + 17 = 0$}
		{Cho điểm $M\left(-\dfrac{3}{2};0;0\right)$ thuộc mặt phẳng $(P)$. Gọi đường thẳng $\Delta$ là hình  chiếu của đường thẳng $BM$ lên $(P)$. Đường thẳng $\Delta$ có một vectơ chỉ phương là $\overrightarrow{u}_{\Delta} = (68;44;10)$}
	 \loigiai{
		\begin{itemchoice}
			\itemch Khoảng cách từ  điểm $A$ đến mặt phẳng $(P)$ được tính bằng công thức
			\[d(A,(P)) = \dfrac{|2(-3) - 4.1 + 4.2 + 3|}{\sqrt{2^2  + (-4)^2 + 4^2}} = \dfrac{1}{6}.\]
			\itemch Vectơ pháp tuyến của mặt phẳng $(P)$ là $\overrightarrow{n} = (1;-2;2)$. \\ Do đó, đường thẳng  $d$ vuông góc với $(P)$ nên có vectơ chỉ phương là $\overrightarrow{n}$. Phương trình tham số của $d$ là $\heva{&x = -3 + t \\ &y = 1 - 2t \\ &z = 2 + 2t}$.
			 \itemch Mặt phẳng $(\alpha)$ đi qua hai điểm $A$, $B$ đồng thời vuông góc với $(P)$ nên có cặp vectơ chỉ phương $\overrightarrow{n} = (1;-2;2)$ và $\overrightarrow{AB} = (4;4;-4)$. \\
			Khi đó, vectơ pháp  tuyến của $(\alpha)$ là $\overrightarrow{n}_{\alpha} = \left[\overrightarrow{n}, \overrightarrow{AB} \right] = (0;3;3)$.\\
			Phương trình tổng quát của $(\alpha)$ là $0x + 3y + 3z + D = 0 \Leftrightarrow  y + z - 3 = 0$.
			\itemch Có $\overrightarrow{MB} = (2;5;-2)$ và $\overrightarrow{b} =  2\overrightarrow{MB} = (5;10;-4)$. \\
			Đặt $\overrightarrow{c}  = \left[\overrightarrow{b}, \overrightarrow{n}\right] = (12;-14;-20)$. \\
			Khi đó, $\overrightarrow{u}_{\Delta} = \left[\overrightarrow{n},  \overrightarrow{c}\right] = (68;44;10)$.
		\end{itemchoice}
	} 
\end{ex}

% nguồn: Phan1_Lop_12_Hinh_hoc_Chuong_5
\begin{ex}%[2H5H2-3]%[ma:7d709dff]
	Trong không gian  $Oxyz$, cho mặt phẳng $(P)\colon x-2y+3=0$ và điểm $A(2; 1; 2)$.
	 \choiceTF
		{\True  Khoảng cách từ điểm $A$ đến $(P)$ bằng $\dfrac{3}{\sqrt{5}}$}
		{\True Đường thẳng $d$ đi qua điểm $A$ và vuông góc với mặt phẳng $(P)$ có phương trình tham  số là $\heva{&x=2+t \\
		&y=1-2t  \\&z=2}$}
		{Mặt phẳng $(Q)$  qua $2$ điểm $B(1;1;0)$, $C(-2;1;1)$ và vuông góc với mặt phẳng $(P)$ có phương trình $(Q)\colon 2x+y=0$}
		{Đường thẳng $d'$ là giao tuyến  của hai mặt phẳng $(P)$ và $(Q)$ có (d) phương trình tham số là $\heva{&x=-3+4t \\
		 &y=2t \\&z=\dfrac{3}{2}-t}$}
	\loigiai{ 
		\begin{itemchoice}
			 \itemch Ta có $\mathrm{d}\left(A,(P)\right)=\dfrac{|2-2.1+3|}{\sqrt{1+(-2)^2}}=\dfrac{3}{\sqrt{5}}$.
			\itemch Đường thẳng $d$ đi qua điểm $A(2; 1; 2)$ và có vectơ chỉ phương $\vec{u}=(1;-2; 0)$ nên có phương trình là  $\heva{&x=2+t \\& y=1-2t \\ &z=2.}$
			\itemch Ta có  $(Q)$ chứa $B$, $C$ nên $\overrightarrow{BC}=(-3; 0;1)$ là một vectơ chỉ phương của $(Q)$.\\
			Lại có $(Q)$ vuông góc với  $(P)$ nên $\vec{n}_P=(1;-2;0)$ là một vectơ chỉ phương của $(Q)$.\\
			Khi đó $\left[\overrightarrow{BC}, \vec{n}_P\right]=(2;1;6)$ là một  vectơ pháp tuyến của $(Q)$.\\
			Vậy $(Q)\colon 2x+y+6z-3=0$. 
			\itemch Ta có $(P) \cap(Q)=d'$ nên $d'$ có dạng $\heva{&x-2y+3=0\\& 2x+y+6z-3=0}$ đi qua điểm $M\left(-3; 0;  \dfrac{3}{2}\right)$\\
			Gọi  $\vec{u}$ là một vectơ chỉ phương của đường thẳng $d'$, ta có $\vec{u}=\left[\vec{n}_P, \vec{n}_Q\right]$ với $\vec{n}_P=(1;-2; 0),\\ \vec{n}_Q=(2;1;6)$ nên $\vec{u}=(-12;-6; 5)$.\\
			Vậy phương trình $d'$ là $\heva{&x=-3-12t \\& y=-6t \\&  z=\dfrac{3}{2}+5t.}$
		\end{itemchoice}
	 }
\end{ex}

% nguồn: Phan1_Lop_12_Hinh_hoc_Chuong_5
\begin{ex}%[2H5H2-3]%[ma:75ca7b24]
	Trong không gian $Oxyz$, cho hai điểm $A(0;2;-1)$,  $B(2;2;-1)$ và đường thẳng\\
	$\Delta \colon \heva{&x=1+2t\\& y=3+t\\& z=4-5t}  (t\in\mathbb{R}).$
	 \choiceTF
		{Một vectơ chỉ  phương của đường thẳng $\Delta$ là $\vec{u}=(1;3;4)$}
		{Phương trình mặt cầu đường kính $AB$ là  $(x-1)^2+(y+2)^2+(z-1)^2=2$}
		{\True Phương  trình chính tắc của đường thẳng $d$ đi qua điểm $A$ và song song với đường thẳng $\Delta$ là $\dfrac{x}{2}=\dfrac{y+2}{1}=\dfrac{z-1}{-5}$}
		{Gọi $\alpha$ là góc  giữa hai đường thẳng $\Delta$ và $AB$, khi đó $\cos\alpha=\dfrac{\sqrt{30}}{15}$}
	 \loigiai{
		\begin{itemchoice}
			\itemch Vectơ chỉ phương của $\Delta$ là $\vec{u}=(2;1;-5)$. 
			\itemch Gọi $I(1;2;-1)$ là trung điểm $AB$. Khi đó  $\left|\overrightarrow{IA}\right|=\sqrt{(0-1)^2+(2-2)^2+(-1+1)^2}=1$.\\
			Phương trình  mặt cầu đường kính $AB$ là $(x-1)^2+(y-2)^2+(z+1)^2=1$.
			\itemch  Phương trình chính tắc của đường thẳng $d$ có dạng $\dfrac{x}{2}=\dfrac{y-2}{2}=\dfrac{z+1}{-5}$.
			\itemch $\cos\left(\Delta,  AB\right)=\cos\left(\overrightarrow{AB},\vec{u}_{\Delta}\right)=\dfrac{|2\cdot2+0\cdot1+0\cdot(-5)|}{\sqrt{2^2+0^2+0^2}\cdot\sqrt{2^2+1^2+(-5)^2}}=\dfrac{\sqrt{30}}{30}$.
		\end{itemchoice}
	} 
\end{ex}

% nguồn: Phan1_Lop_12_Hinh_hoc_Chuong_5
\begin{ex}%[2H5V2-3]%[ma:4f102372]
	Trong không gian với hệ tọa độ $Oxyz$,  cho đường thẳng $d\colon\dfrac{x+4}{3}=\dfrac{y-5}{-4}=\dfrac{z+2}{1}$.
	\choiceTF 
		 {Điểm $M(3;-4;1)$ nằm trên đường thẳng $d$}
		{\True Một  vectơ chỉ phương của $d$ là $\vec{u}=(-3;4;-1)$}
		{Đường thẳng $d$ song song với  đường thẳng $\Delta\colon\heva{&x=-4+3t\\& y=5-4t\\& z=t}$}
		{\True Đường thẳng $\Delta_1$ song song với $d$ và cắt  cả hai đường thẳng $d_1\colon\dfrac{x-1}{3}=\dfrac{y+1}{1}=\dfrac{z-2}{2}$ và $d_2\colon\dfrac{x+2}{2}=\dfrac{y-3}{4}=\dfrac{z}{1}$ có phương trình là $\Delta_1\colon\dfrac{x+4}{3}=\dfrac{y+1}{-4}=\dfrac{z+1}{1}$}
	 \loigiai{
		\begin{itemchoice}
			\itemch Thay tọa độ điểm $M(3;-4;1)$ vào phương trình $d$,  ta được
			$\dfrac{3+4}{3}=\dfrac{-4-5}{-4}=\dfrac{1+2}{1}$ (sai) nên điểm $M(3;-4;1)$ không nằm  trên đường thẳng $d$.
			\itemch Vectơ chỉ phương của $d$ là $\vec{v}=(3;-4;1)$.  \\
			Ta có $\vec{u}=(-3;4;-1)=-(3;-4;1)=-\vec{v}$. Vì $\vec{u}$ cùng phương với $\vec{v}$, nên $\vec{u}$ cũng là một vectơ chỉ phương của  $d$.
			\itemch Đường thẳng $d$ có một vectơ chỉ phương $\vec{u}_d=(3;-4;1)$ và đi  qua $A(-4;5;-2)$.\\
			Đường thẳng $\Delta$  có một vectơ chỉ phương $\vec{u}_{\Delta}=(3;-4;1)$ và đi qua $B(-4;5;0)$.\\
			Vì $\vec{u}_d=\vec{u}_{\Delta}$ nên $d$ và $\Delta$ song song hoặc trùng  nhau.\\
			Thay tọa độ điểm $M(-4;5;0)$ vào phương trình $d$, ta  được
			$\dfrac{-4+4}{3}=\dfrac{5-5}{-4}=\dfrac{0+2}{1}$  (sai).\\
			Vậy  $B$ không thuộc $d$. Do đó, đường thẳng $d$ song song với đường thẳng $\Delta$.
			\itemch Đường thẳng $\Delta_1$ song  song với $d$ nên $\Delta_1$ có vectơ chỉ phương $\vec{u}_{\Delta_1}=(3;-4;1)$.\\
			Gọi $M$, $N$ là  giao điểm của $\Delta_1$ và $d_1$, $d_2$.\\
			Khi đó $M\in d_1\Rightarrow M(1+3t;-1+t; 2+2t)$ và $N\in d_2\Rightarrow N(-2+2t';  3+4t'; t')$.\\
			Vectơ chỉ phương của $\Delta_1$  là $\vec{MN}=(-3+2t'-3t;4+4t'-t;-2+t'-2t)$.\\
			Vì $\Delta_1$  song song với $d$, nên $\vec{MN}$ cùng phương với $\vec{u}_d=(3;-4;1)$ hay
			$$\dfrac{-3+2t'-3t}{3}=\dfrac{4+4t'-t}{-4}=\dfrac{-2+t'-2t}{1}.$$ 
			 Từ đó suy ra $\heva{&3t-t'=-3\\&9t-8t'=4}\Rightarrow \heva{&t=\dfrac{4}{3}\\&t'=-1.}$\\
			Suy ra điểm $N(-4;-1;-1)$  nằm trên $\Delta_1$.\\
			Vậy $\Delta_1$ đi qua $N(-4;-1;-1)$ và có vectơ chỉ phương $\vec{u}_{\Delta_1}=(3;-4;1)$  có phương trình là $ \dfrac{x+4}{3}=\dfrac{y+1}{-4}=\dfrac{z+1}{1}$.
		\end{itemchoice}
	 }
\end{ex}

% nguồn: Phan1_Lop_12_Hinh_hoc_Chuong_5
\begin{ex}%[2H5V2-3]%[ma:d9416841]
	Trong  không gian $Oxyz$, cho điểm $M(2;1;0)$ và đường thẳng\\ $d\colon \dfrac{x-1}{2}=\dfrac{y+1}{1}=\dfrac{z}{-1}$.
	\choiceTF 
		{\True Đường thẳng $d$ có một véc-tơ  chỉ phương là $\overrightarrow{u}=(2;1;-1)$}
		{\True Mặt phẳng $(P)$ đi qua điểm $M$ và vuông góc với  $d$ có phương trình tổng quát là $2x+by+cz+d=0$. Khi đó $b+c+d=-5$}
		{Gọi $M'$ là điểm đối xứng với $M$ qua $d$. Khi đó  $M'(1;0;-2)$}
		{\True Phương trình đường thẳng $\Delta$ đi qua điểm $M$ cắt và vuông góc với đường thẳng $d$ $\dfrac{x-2}{1}=\dfrac{y-1}{a}=\dfrac{z}{b}$. Khi  đó $a+b=-6$}
	\loigiai{ 
		\begin{itemchoice}
			\itemch Véc-tơ chỉ  phương của $d$ là $\overrightarrow{u}=(2;1;-1)$.
			\itemch Phương trình mặt phẳng vuông góc với $d$ có pháp tuyến là $\overrightarrow{n}=\overrightarrow{u}=(2;1;-1)$. \\ 
			Phương trình tổng quát của mặt phẳng $(P)$ là $2x+y-z+d=0$.\\ 
			Thay điểm $M$ vào $(P)$ ta được $2\cdot 2+1-1 \cdot 0+d=0 \Rightarrow 4+1+d=0 \Rightarrow  d=-5$.\\
			Do đó,  $b+c+d=1-1-5=-5$.
			 \itemch Phương trình mặt phẳng $(P)\colon 2x+y-z-5=0$.\\
			Đường thẳng $d$  có phương trình tham số
			 $\heva{&x=1+2t\\&y=-1+t\\&z=-t}$\\
			Gọi  $H$ là giao điểm của $d$ với mặt phẳng $(P)$, khi đó $H=(1+2t,-1+t,-t)$.\\
			Thay  vào phương trình mặt phẳng $(P)$ ta được\\
			 $2(1+2t)+(-1+t)-(-t)-5=0 \Rightarrow t=\dfrac{2}{3}$.\\
			Suy ra tọa độ  điểm $H$ là
			 $H=\left(\dfrac{7}{3},-\dfrac{1}{3},-\dfrac{2}{3}\right)$.\\
			Vì $H$ là trung điểm của đoạn $MM'$ nên $M'\left( \dfrac{8}{3};\dfrac{-5}{3};\dfrac{-4}{3}\right) $. 
			\itemch Gọi $N$ là giao điểm  của $\Delta$ và $d$. \\
			Vì  $N$ thuộc $d$, nên tọa độ $N=(1+2t,-1+t,-t)$.\\
			 $\overrightarrow{MN}=(1+2t-2,-1+t-1,-t-0)=(-1+2t,-2+t,-t)$.\\
			 Đường thẳng $\Delta$ vuông góc với $d$ nên $\overrightarrow{MN} \perp \overrightarrow{u}\Rightarrow \overrightarrow{MN} \cdot \overrightarrow{u}=0$.\\
			$\Rightarrow(-1+2t)\cdot 2+(-2+t)\cdot 1+(-t)\cdot (-1)=0$  \\
			$\Rightarrow  t=\dfrac{2}{3}$\\
			Thay $t=\dfrac{2}{3}$ vào $\overrightarrow{MN}$,  ta được\\
			$\overrightarrow{MN}=\left(\dfrac{1}{3},-\dfrac{4}{3},-\dfrac{2}{3}\right)=\dfrac{1}{3}\cdot(1,-4,-2)$.\\ 
			Chọn véc-tơ chỉ phương của $\Delta$  là $\overrightarrow{v}=(1,-4,-2)$.\\
			Phương trình chính tắc của  $\Delta$ là
			 $\dfrac{x-2}{1}=\dfrac{y-1}{-4}=\dfrac{z}{-2}$.\\
			Vậy $a=-4$, $b=-2$  $\Rightarrow a+b=-6$.
		\end{itemchoice}
	 }
\end{ex}

% nguồn: Phan1_Lop_12_Hinh_hoc_Chuong_5
\begin{ex}%[2H5V2-3]%[ma:59cef285]
	Trong không gian với hệ tọa độ $Oxyz$, cho điểm $M(2;1;0)$ và đường thẳng  \break{$d\colon\dfrac{x-1}{2}=\dfrac{y+1}{1}=\dfrac{z}{-1}$}.
	\choiceTF 
		{Điểm $M(1;-1;1)$ nằm  trên đường thẳng $d$}
		{Một vectơ chỉ phương  của $d$ là $\vec{u}=(2;1;-1)$}
		{\True Đường thẳng  $d$ song song với đường thẳng $\Delta\colon\heva{&x=1+2l\\ &y=-1+l\\ &z=-l}$}
		{\True  Gọi $\Delta_1$ là đường thẳng đi qua $M$, cắt và vuông góc với $d$ thì phương trình đường thẳng $\Delta_1$ là $\heva{&x=2+t\\ &y=1-4t\\ &z=-2t}$}
	 \loigiai{
		\begin{itemchoice}
			\itemch Thay tọa độ điểm  $M(1;-1;1)$ vào phương trình $d$, ta được
			$\dfrac{1-1}{2}=\dfrac{-1+1}{1}=\dfrac{1}{-1}$ (sai) nên điểm $M(1;-1;1)$ không nằm trên đường thẳng $d$. 
			\itemch Từ phương trình chính tắc của $d$, suy ra một vectơ  chỉ phương của $d$ là $\vec{u}=(2;1;-1)$.
			\itemch Đường thẳng $d$ có một vectơ chỉ phương  $\vec{u}=(2;1;-1)$ và đi qua $A(1;-1;0)$.\\
			Đường thẳng $\Delta\colon\heva{&x=1+2l\\ &y=-1+l\\ &z=-l}$ có một  vectơ chỉ phương $\vec{u}_{\Delta}=(2;1;-1)$ và đi qua $B(1;-1;0)$.\\
			Vì $\left[\vec{u},\vec{u}_{\Delta}\right]=\vec{0}$ và điểm $A(1;-1;0)$ của $d$ cũng thuộc $\Delta$, nên đường thẳng $d$ trùng với đường thẳng  $\Delta$.
			\itemch  Gọi $H$ là hình chiếu vuông góc của $M$ lên $d$. Đường thẳng $\Delta_1$ cần tìm là đường thẳng $MH$.\\
			Phương trình tham số của đường thẳng $d$ là $\heva{&x=1+2t'\\ &y=-1+t'\\  &z=-t'}$.\\
			Vì $H\in d$ nên $H(1+2t';-1+t';-t')$ và  $\vec{MH}=(2t'-1;t'-2;-t')$.\\
			 Mà $\vec{MH}\cdot\vec{u}_d=0$ nên
			 $2(2t'-1)+1(t'-2)+(-1)(-t')=0\Leftrightarrow 6t'-4=0\Leftrightarrow t'=\dfrac{2}{3}$.\\
			 Khi đó, $\vec{MH}=\left(\dfrac{1}{3};-\dfrac{4}{3};-\dfrac{2}{3}\right) $.\\
			Ta có thể chọn vectơ chỉ phương cho  $\Delta_1$ là $\vec{u}_1=3\vec{MH}=(1;-4;-2)$.\\
			Đường thẳng $\Delta_1$ đi qua $M(2;1;0)$ và  có vectơ chỉ phương $\vec{u}_1=(1;-4;-2)$ nên có phương trình là $\heva{&x=2+t\\ &y=1-4t\\ &z=-2t}$.
		\end{itemchoice}
	} 
\end{ex}


%%%%%%%%%%%%%%%% TRẢ LỜI NGẮN %%%%%%%%%%%%%%%%

% nguồn: Phan2_Lop_12_Hinh_hoc_Chuong_5
\begin{bt}%[2H5H2-3]%[ma:f5e102de]
	Trong  không gian $Oxyz$, cho hai đường thẳng $d_1\colon \heva{&x=t \\
	&y=-1-4t\\&z=6+6t}$ và $d_2\colon \dfrac{x}{2}=\dfrac{y-1}{1}=\dfrac{z+2}{-5}$. Phương trình chính tắc của đường thẳng $d_3$ qua $M(1;-1; 2)$ và vuông góc với cả $d_1$, $d_2$ có dạng $\dfrac{x-1}{14}=\dfrac{y+1}{b}=\dfrac{z-a}{c}$.  Tính $a+b+c$.
	\shortans{28} 
	 \loigiai{
		$d_1$ có một véc-tơ chỉ phương là $\overrightarrow{u}_1=(1;-4;6)$.\\ 
		$d_2$ có một véc-tơ  chỉ phương là $\overrightarrow{u}_2=(2;1;-5)$.\\
		Vì $d_3$ vuông góc với cả $d_1$, $d_2$ nên $d_3 $ có một  véc-tơ chỉ phương là $\overrightarrow{u}_3=\left[\overrightarrow{u}_1,\overrightarrow{u}_2\right]=(14;17;9)$.\\
		Suy ra  $d_3\colon \dfrac{x-1}{14}=\dfrac{y+1}{17}=\dfrac{z-2}{9}$.\\
		Do đó  $a=2$, $b=17$, $c=9$.\\
		Vậy  $a+b+c=28$.
	 }
\end{bt}

% nguồn: Phan2_Lop_12_Hinh_hoc_Chuong_5
\begin{bt}%[2H5H2-3]%[ma:c1536f28]
	Trong không gian $Oxyz$, cho điểm $A(2;1;3)$, đường thẳng $d:\left\{\begin{aligned}
	x&=-1+t  \\
	y&=1-2t \\ 
	 z&=2+t
\end{aligned}\right. \; (t \in \mathbb{R})$ và mặt phẳng $(P): x+y - 2z+2=0$. Phương trình chính tắc của đường thẳng $\Delta$ đi qua $A$, song song với mặt phẳng $(P)$ và vuông góc với đường thẳng $d$ có dạng $\dfrac{x+a}{b}=\dfrac{y-1}{c}=\dfrac{z+d}{3}$. Giá trị của biểu thức $T=a+b+c+d$ bằng bao nhiêu?
\shortans{1} 
 \loigiai{
	Đường thẳng $d$ có một vectơ chỉ phương  $\vec{u}=(1;-2;1)$. \\
	Mặt phẳng $(P)$ có một vectơ pháp tuyến  $\vec{n}=(1;1;-2)$. \\
	Ta có  $\left[\vec{u},\vec{n}\right]=(3;3;3)$. \\
	Vì đường thẳng $\Delta$ song song với $(P)$ và vuông góc  với $d$ nên có một vectơ chỉ phương \\
	 $$\vec{a}=\left[\vec{u},\vec{n}\right]=(3;3;3).$$
	Đường thẳng $\Delta$ đi qua điểm $A(2;1;3)$ và có một vectơ chỉ phương  $\vec{a}=(3;3;3)$ nên có phương trình chính tắc \\
	$$\dfrac{x-2}{3}=\dfrac{y-1}{3}=\dfrac{z-3}{3}.$$ 
	Do đó $a=-2$, $b=3$, $c=3$ và $d=-3$.  \\
	 Vậy $T=a+b+c+d=1$.
 }
\end{bt}

% nguồn: Phan2_Lop_12_Hinh_hoc_Chuong_5
\begin{bt}%[2H5H2-3]%[ma:6c456fbc]
	Trong  không gian $Oxyz$, cho $A(1; 3; -2)$, $B(1;1;5)$. Phương trình tham số của đường thẳng đi qua hai điểm $A$, $B$ có dạng $\heva{& x=a \\ & y=b-2t \\ & z=-2+ct}$. Tính $a+b+c$.
	 \shortans{11}
	\loigiai{Ta có $\vec{AB}  = (0; -2; 7)$.\\
		Đường  thẳng $AB$ đi qua $A(1; 3; -2)$ và nhận $\vec{AB} = (0;-2;7)$ làm vectơ chỉ phương có phương trình là $\heva{& x=1 \\ & y=3-2t \\ & z=-2+7t }$, $ t \in \mathbb{R}$.\\
		Suy ra $a =  1$, $b = 3$, $c = 7$. Khi đó $a + b + c = 11$.
	 }
\end{bt}

% nguồn: Phan2_Lop_12_Hinh_hoc_Chuong_5
\begin{bt}%[2H5H2-3]%[ma:4974bea3]
	Trong không gian $Oxyz$, cho ba điểm $A(1;0;1)$, $B(1;1;0)$ và $C(3;4;-1)$. Phương trình đường thẳng đi qua $A$ và song song với $B C$ có dạng $\dfrac{x-1}{2}=\dfrac{y}{m}=\dfrac{z-1}{n}$ $(m$; $n  \in \mathbb{Z})$. Tính giá trị $S=m+n$
	 \shortans{2}
	\loigiai{Đường thẳng $d$ đi qua $A$ và song song với $BC$ nhận $\vec{BC}=(2;3;-1)$ làm một véc tơ  chỉ phương.\\
		Phương trình  của đường thẳng $d\colon \dfrac{x-1}{2}=\dfrac{y}{3}=\dfrac{z-1}{-1}$.\\
		 Khi đó $\heva{&m=3 \\& n=-1} \Rightarrow S=m+n=2.$
	 }
\end{bt}

% nguồn: Phan2_Lop_12_Hinh_hoc_Chuong_5
\begin{bt}%[2H5H2-3]%[ma:5ea3bf9f]
	Trong không gian với hệ trục tọa độ $Oxyz$, cho hai điểm $A(1; 2; -1)$, $B(2; -1; 3)$. đường thẳng đi qua $2$ điểm $A$, $B$ có phương trình dạng $\Delta \colon \heva{& x = 2t \\ & y = y_0+ bt \\ & z= z_0 +  ct}$. Tính $y_0 + z_0 + b + c$.
	 \shortans{2}
	\loigiai{ 
		Đường thẳng $\Delta$ có  vectơ chỉ phương là $\vec{AB} = (1; -3; 4) = \dfrac{1}{2}(2; -6; 8) \Rightarrow b=-6, c=8$.\\
		Thay tọa độ điểm $B$ vào phương  trình $\Delta$ ta được $\heva{& 2 = 2t \\ & -1 = y_0 - 6t \\ & 3 = z_0 + 8t} \Leftrightarrow \heva{& t=1 \\ & y_0 = 5 \\ & {z_0}=-5.}$\\
		Từ  đó ta có $y_0 + z_0 + b + c = 5 - 5 - 6 + 8 = 2$.
	} 
\end{bt}

% nguồn: Phan2_Lop_12_Hinh_hoc_Chuong_5
\begin{bt}%[2H5V2-3]%[ma:c119ac92]
	Trong không gian $Oxyz$, cho điểm $A(1;2;-3)$, đường thẳng $\Delta \colon \heva{&x=1-2t \\ &y=5-t \\ &z=-2+3t }$ và mặt phẳng $(P)\colon 3x-5y+z-1=0$. Một đường thẳng $d$ đi qua điểm $A$, $d$ cắt đường thẳng $\Delta$  và $d$ song song với mặt phẳng $(P)$ có một vectơ chỉ phương là $\vec{u}=(a;b;1)$. Tính giá trị biểu thức $T=2a+5b$ (làm tròn kết quả đến hàng phần chục).\shortans{-2,2}
	\loigiai{ 
		Gọi $M$ là giao điểm của đường thẳng  $d$ và đường thẳng $\Delta$.\\
		Vì $M \in \Delta$, tọa độ của $M$ có dạng $M(1-2t; 5-t; -2+3t)$. 
		Đường thẳng $d$ đi qua $A(1;2;-3)$ và  $M(1-2t; 5-t; -2+3t)$.\\
		Một vectơ chỉ phương của $d$ là $\overrightarrow{AM}  = (-2t; 3-t; 1+3t)$.\\
		Đường thẳng $d$  song song với mặt phẳng $(P)$, nên vectơ chỉ phương $\overrightarrow{AM}$ của $d$ vuông góc với vectơ pháp tuyến $\vec{n}_P = (3;-5;1)$ của $(P)$.\\
		Do đó, $\overrightarrow{AM} \cdot \vec{n}_P  = 0\Leftrightarrow (-2t)\cdot3 + (3-t)\cdot(-5) + (1+3t)\cdot1 = 0\Leftrightarrow t = 7$.\\
		Với $t=7$, vectơ chỉ phương của $d$ là $\overrightarrow{AM} = (-14;  -4; 22)$.\\
		Để vectơ chỉ phương này có cao độ bằng $1$, ta chia  tất cả các thành phần cho $22$:
		 $\vec{u}' = \left(\dfrac{-14}{22}; \dfrac{-4}{22}; \dfrac{22}{22}\right) = \left(-\dfrac{7}{11}; -\dfrac{2}{11}; 1\right)$.\\
		So sánh với $\vec{u}(a;b;1)$, ta có $a = -\dfrac{7}{11}$ và  $b = -\dfrac{2}{11}$.\\
		Giá trị biểu thức  $T = 2a + 5b= -\dfrac{24}{11}\approx -2{,}2$.
	} 
\end{bt}


%%%%%%%%%%%%%%%% TỰ LUẬN %%%%%%%%%%%%%%%%

% nguồn: Phan2_Lop_12_Hinh_hoc_Chuong_5
\begin{bt}%[2H5H2-3]%[ma:376f1c36]
	Trong không gian với hệ toạ độ $Oxyz$,  viết phương trình tham số và phương trình chính tắc của đường thẳng $\Delta$ trong mỗi trường hợp sau
	\begin{itemize}
		\item[a)] $\Delta$ đi qua điểm $A(-1;3;2)$ và có vectơ chỉ  phương $\overrightarrow{u}=(-2;3;4)$;
		\item[b)] $\Delta$  đi qua hai điểm $M(2;-1;3)$ và $N(3;0;4)$.
	\end{itemize}
	 \loigiai{
		\begin{itemize}
			\item  [a)] Phương trình tham số của đường thẳng $\Delta$ đi qua điểm $A(-1;3;2)$ và có vectơ chỉ phương $\overrightarrow{u}=(-2;3;4)$ là
			 $$\heva{
			&x=-1-2t\\ 
			&y=3+3t  \\
			&z=2+4t 
			} \quad (t \text{ là tham  số}).$$
			Phương trình chính tắc của đường thẳng $\Delta$ đi qua điểm $A(-1;3;2)$ và có vectơ chỉ phương $\overrightarrow{u}=(-2;3;4)$  là
			 \[
			 \dfrac{x+1}{-2}=\dfrac{y-3}{3}=\dfrac{z-2}{4}.
			 \]
			\item [b)] Ta có  $\overrightarrow{MN}=(1;1;1)$ là một vectơ chỉ phương của đường thẳng $\Delta$. \\
			Phương trình tham số của đường  thẳng $\Delta$ là
			$$\heva{&x=2+t' \\ &y=-1+t' \\ &z=3+t'} \quad (t' \text{ là  tham số}).$$
			Phương trình chính tắc của đường  thẳng $\Delta$ là
			$$\dfrac{x-2}{1}=\dfrac{y+1}{1}=\dfrac{z-3}{1}.$$ 
		\end{itemize}
	 }
\end{bt}

% nguồn: Phan2_Lop_12_Hinh_hoc_Chuong_5
\begin{bt}%[2H5H2-3]%[ma:953d78a2]
	Trong không gian với hệ toạ độ $Oxyz$, viết các phương trình tham số và chính tắc của đường thẳng $\Delta$ đi qua hai điểm $A(2;3;-1)$ và  $B(1;-2;4)$.
	 \loigiai{
		Ta có  $\overrightarrow{AB}=(-1;-5;5)$. \\
		Vì $\Delta$ đi qua  hai điểm $A$ và $B$ nên vectơ chỉ phương của $\Delta$ là $\overrightarrow{u}=\overrightarrow{AB}=(-1;-5;5)$. \\
		Đường thẳng $\Delta$ đi qua điểm $A(2;3;-1)$ và có vectơ chỉ phương là  $\overrightarrow{u}=(-1;-5;5)$. \\
		Phương trình tham số của $\Delta$ là 
		$\heva{&x=2-t \\ &y=3-5t \\  &z=-1+5t} \quad (t \in \mathbb{R})$. \\
		Phương trình chính  tắc của $\Delta$ là
		$\dfrac{x-2}{-1}=\dfrac{y-3}{-5}=\dfrac{z+1}{5}$. 
	 }
\end{bt}

% nguồn: Phan2_Lop_12_Hinh_hoc_Chuong_5
\begin{bt}%[2H5H2-3]%[ma:c0a4e519]
	Trong không gian với hệ toạ độ $Oxyz$, một viên đạn được bắn ra từ điểm $A(1;3;4)$ và trong $3$ giây, đầu đạn đi với vận tốc không đổi, vectơ vận tốc (trên giây) là $\overrightarrow{v}=(2;1;6)$. Hỏi viên đạn trên có bắn trúng mục tiêu trong mỗi tình huống  sau hay không?
	\begin{itemize}
		\item[a)] Mục tiêu  đặt tại điểm $M\left(7;\dfrac{7}{2};21\right)$;
		\item[b)] Mục tiêu đặt  tại điểm $N(-3;1;-8)$.
	\end{itemize}
	 \loigiai{
		Phương trình mô  tả quỹ đạo chuyển động của viên đạn là
		$\heva{&x=1+2t \\  &y=3+t \\ &z=4+6t}$.
		\begin{itemize}
			\item[a)] Thay tọa độ điểm $M$ vào phương trình chuyển  động ta có
			$$\heva{&7=1+2t \\ &\dfrac{7}{2}=3+t  \\ &21=4+6t}
			\Leftrightarrow \heva{&t=3 \\ &t=\dfrac{1}{2} \\ &t=\dfrac{17}{6}.}$$ 
			Ta thấy các giá trị $t$ đều khác nhau nên điểm $M$ không nằm trên quỹ đạo chuyển động của  viên đạn. Vậy viên đạn không bắn trúng mục tiêu đặt tại $M$.
			\item[b)] Thay tọa độ điểm $N$ vào phương trình chuyển động ta có 
			$$\heva{&-3=1+2t  \\ &1=3+t \\ &-8=4+6t} \\
			 \Leftrightarrow t=-2.$$
			Suy ra điểm $N$ nằm trên quỹ đạo chuyển động của viên đạn. Do đó viên đạn  bắn trúng mục tiêu tại $N$.
		\end{itemize}
	} 
\end{bt}

% nguồn: Phan2_Lop_12_Hinh_hoc_Chuong_5
\begin{bt}%[2H5V2-3]%[ma:e1722313]
	Trong không gian với hệ toạ độ $Oxyz$, cho hai mặt phẳng $(\alpha) \colon 2x+y-z+3=0$ và $(\beta) \colon  x+y+z-1=0$. Viết phương trình chính tắc đường thẳng $\Delta$ là giao tuyến của hai mặt phẳng $(\alpha)$ và $(\beta)$.
	 \loigiai{
		Mặt phẳng $(\alpha)$ có một vectơ pháp tuyến $\overrightarrow{n}_{\alpha}=(2;1;-1)$.  \\
		Mặt phẳng $(\beta)$ có một vectơ pháp tuyến $\overrightarrow{n}_{\beta}=(1;1;1)$.  \\
		Suy ra một vectơ chỉ phương của giao  tuyến $\Delta$ là $$\overrightarrow{u}_\Delta=\left[\overrightarrow{n}_{\alpha}, \overrightarrow{n}_{\beta}\right]=(2;-3;1).$$
		Gọi $M$ là một  điểm trên $\Delta$. \\
		Suy ra $M$ là một  điểm chung của hai mặt phẳng $(\alpha)$ và $(\beta)$. \\
		Khi đó tọa độ điểm $M$ thỏa  mãn hệ phương trình
		$$\heva{&2x+y-z+3=0 \\  &x+y+z-1=0.}$$
		Lấy  $x=0$, ta có $$\heva{&y-z+3=0 \\ &y+z-1=0} \Leftrightarrow \heva{&y=-1 \\ &z=2.}$$
		Do  đó $M(0;-1;2)$. \\
		Vậy phương trình chính tắc của  đường thẳng $\Delta$ là $$\dfrac{x}{2}=\dfrac{y+1}{-3}=\dfrac{z-2}{1}.$$
	} 
\end{bt}

% nguồn: Phan2_Lop_12_Hinh_hoc_Chuong_5
\begin{bt}%[2H5V2-3]%[ma:25be49fd]
	Trong không gian với  hệ trục toạ độ $Oxyz$, cho mặt phẳng $(P) \colon x+y-5z+4=0$ và đường thẳng $d \colon \dfrac{x+1}{2}=\dfrac{y+1}{1}=\dfrac{z+5}{6}$. Gọi $(Q)$ là mặt phẳng chứa $d$ và vuông góc với $(P)$. Xác định phương trình giao tuyến $d'$ của $(Q)$ và $(P)$?
	 \loigiai{
		Gọi đường thẳng $d'$  là hình chiếu vuông góc của đường thẳng $d$ trên $(P)$. \\
		 Đường thẳng $d$ đi qua điểm $A(-1;-1;-5)$ và có $\overrightarrow{u}_d=(2;1;6)$. \\
		 Mặt phẳng $(P)$ có vectơ pháp tuyến $\overrightarrow{n}_P=(1;1;-5)$. \\
		Mặt phẳng $(Q)$ chứa $d$ và vuông góc với $(P)$ nên một vectơ pháp tuyến của $(Q)$ là $\overrightarrow{n}_Q=[\overrightarrow{n}_P,  \overrightarrow{u}_d]=(11;-16;-1)$. \\
		 Phương trình của mặt phẳng $(Q)$ là $11x-16y-z-10=0$. \\
		Do $(P) \cap (Q)=d'$ nên một vectơ chỉ phương  của đường thẳng $d'$ là
		 $$\left[\overrightarrow{n}_P,\overrightarrow{n}_Q\right]=(-81;-54;-27).$$
		Chọn $\overrightarrow{u}_1=-\dfrac{1}{27} \left[\overrightarrow{n}_P,\overrightarrow{n}_Q\right]=(3;2;1)$  là một vectơ chỉ phương của $d'$. \\
		 Phương trình tham số của đường thẳng $d$ là
		$\heva{&x=-1+2t \\ &y=-1+t  \\ &z=-5+6t}$. \\
		 Gọi $I=d \cap (P)$. Khi đó tọa độ $I$ thỏa mãn hệ phương trình
		$\heva{&x+y-5z+4=0 \\ &x=-1+2t \\ &y=-1+t \\ &z=-5+6t}$. \\ 
		Giải hệ ta được  $(x;y;z)=(1;0;1)$. \\
		Đường thẳng $d'$ đi qua $I(1;0;1)$ và có một vectơ chỉ phương  $\overrightarrow{u}_1=(3;2;1)$ nên có phương trình chính tắc
		 $$\dfrac{x-1}{3}=\dfrac{y}{2}=\dfrac{z-1}{1}.$$
	 }
\end{bt}
